%% bare_jrnl.tex
%% V1.4b
%% 2015/08/26
%% by Michael Shell
%% see http://www.michaelshell.org/
%% for current contact information.
%%
%% This is a skeleton file demonstrating the use of IEEEtran.cls
%% (requires IEEEtran.cls version 1.8b or later) with an IEEE
%% journal paper.
%%
%% Support sites:
%% http://www.michaelshell.org/tex/ieeetran/
%% http://www.ctan.org/pkg/ieeetran
%% and
%% http://www.ieee.org/

%%*************************************************************************
%% Legal Notice:
%% This code is offered as-is without any warranty either expressed or
%% implied; without even the implied warranty of MERCHANTABILITY or
%% FITNESS FOR A PARTICULAR PURPOSE! 
%% User assumes all risk.
%% In no event shall the IEEE or any contributor to this code be liable for
%% any damages or losses, including, but not limited to, incidental,
%% consequential, or any other damages, resulting from the use or misuse
%% of any information contained here.
%%
%% All comments are the opinions of their respective authors and are not
%% necessarily endorsed by the IEEE.
%%
%% This work is distributed under the LaTeX Project Public License (LPPL)
%% ( http://www.latex-project.org/ ) version 1.3, and may be freely used,
%% distributed and modified. A copy of the LPPL, version 1.3, is included
%% in the base LaTeX documentation of all distributions of LaTeX released
%% 2003/12/01 or later.
%% Retain all contribution notices and credits.
%% ** Modified files should be clearly indicated as such, including  **
%% ** renaming them and changing author support contact information. **
%%*************************************************************************


% *** Authors should verify (and, if needed, correct) their LaTeX system  ***
% *** with the testflow diagnostic prior to trusting their LaTeX platform ***
% *** with production work. The IEEE's font choices and paper sizes can   ***
% *** trigger bugs that do not appear when using other class files.       ***                          ***
% The testflow support page is at:
% http://www.michaelshell.org/tex/testflow/



\documentclass[journal]{IEEEtran}
%
% If IEEEtran.cls has not been installed into the LaTeX system files,
% manually specify the path to it like:
% \documentclass[journal]{../sty/IEEEtran}
\usepackage{url}
\usepackage{cite}
\usepackage{graphicx}
\usepackage{float}
\usepackage{booktabs} 
\usepackage{multirow} 
\usepackage{tabularx}
\usepackage{booktabs}
\usepackage{amsmath}
\usepackage{amssymb}
 
\usepackage{tikz}
\usetikzlibrary{positioning, arrows.meta, shapes.multipart}
\usepackage{placeins}
% placeins/\FloatBarrier is used at a few points below to prevent queued
% tables/figures from being placed so late that they visually collide with
% unrelated later content -- the exact class of rendering defect flagged in
% the reviewer assessment for Table XII ("No MRDL" row appearing awkwardly
% aligned).
% Some very useful LaTeX packages include:
% (uncomment the ones you want to load)


% *** MISC UTILITY PACKAGES ***
%
%\usepackage{ifpdf}
% Heiko Oberdiek's ifpdf.sty is very useful if you need conditional
% compilation based on whether the output is pdf or dvi.
% usage:
% \ifpdf
%   % pdf code
% \else
%   % dvi code
% \fi
% The latest version of ifpdf.sty can be obtained from:
% http://www.ctan.org/pkg/ifpdf
% Also, note that IEEEtran.cls V1.7 and later provides a builtin
% \ifCLASSINFOpdf conditional that works the same way.
% When switching from latex to pdflatex and vice-versa, the compiler may
% have to be run twice to clear warning/error messages.






% *** CITATION PACKAGES ***
%
%\usepackage{cite}
% cite.sty was written by Donald Arseneau
% V1.6 and later of IEEEtran pre-defines the format of the cite.sty package
% \cite{} output to follow that of the IEEE. Loading the cite package will
% result in citation numbers being automatically sorted and properly
% "compressed/ranged". e.g., [1], [9], [2], [7], [5], [6] without using
% cite.sty will become [1], [2], [5]--[7], [9] using cite.sty. cite.sty's
% \cite will automatically add leading space, if needed. Use cite.sty's
% noadjust option (cite.sty V3.8 and later) if you want to turn this off
% such as if a citation ever needs to be enclosed in parenthesis.
% cite.sty is already installed on most LaTeX systems. Be sure and use
% version 5.0 (2009-03-20) and later if using hyperref.sty.
% The latest version can be obtained at:
% http://www.ctan.org/pkg/cite
% The documentation is contained in the cite.sty file itself.






% *** GRAPHICS RELATED PACKAGES ***
%
\ifCLASSINFOpdf
  % \usepackage[pdftex]{graphicx}
  % declare the path(s) where your graphic files are
  % \graphicspath{{../pdf/}{../jpeg/}}
  % and their extensions so you won't have to specify these with
  % every instance of \includegraphics
  % \DeclareGraphicsExtensions{.pdf,.jpeg,.png}
\else
  % or other class option (dvipsone, dvipdf, if not using dvips). graphicx
  % will default to the driver specified in the system graphics.cfg if no
  % driver is specified.
  % \usepackage[dvips]{graphicx}
  % declare the path(s) where your graphic files are
  % \graphicspath{{../eps/}}
  % and their extensions so you won't have to specify these with
  % every instance of \includegraphics
  % \DeclareGraphicsExtensions{.eps}
\fi
% graphicx was written by David Carlisle and Sebastian Rahtz. It is
% required if you want graphics, photos, etc. graphicx.sty is already
% installed on most LaTeX systems. The latest version and documentation
% can be obtained at: 
% http://www.ctan.org/pkg/graphicx
% Another good source of documentation is "Using Imported Graphics in
% LaTeX2e" by Keith Reckdahl which can be found at:
% http://www.ctan.org/pkg/epslatex
%
% latex, and pdflatex in dvi mode, support graphics in encapsulated
% postscript (.eps) format. pdflatex in pdf mode supports graphics
% in .pdf, .jpeg, .png and .mps (metapost) formats. Users should ensure
% that all non-photo figures use a vector format (.eps, .pdf, .mps) and
% not a bitmapped formats (.jpeg, .png). The IEEE frowns on bitmapped formats
% which can result in "jaggedy"/blurry rendering of lines and letters as
% well as large increases in file sizes.
%
% You can find documentation about the pdfTeX application at:
% http://www.tug.org/applications/pdftex





% *** MATH PACKAGES ***
%
%\usepackage{amsmath}
% A popular package from the American Mathematical Society that provides
% many useful and powerful commands for dealing with mathematics.
%
% Note that the amsmath package sets \interdisplaylinepenalty to 10000
% thus preventing page breaks from occurring within multiline equations. Use:
%\interdisplaylinepenalty=2500
% after loading amsmath to restore such page breaks as IEEEtran.cls normally
% does. amsmath.sty is already installed on most LaTeX systems. The latest
% version and documentation can be obtained at:
% http://www.ctan.org/pkg/amsmath





% *** SPECIALIZED LIST PACKAGES ***
%
%\usepackage{algorithmic}
% algorithmic.sty was written by Peter Williams and Rogerio Brito.
% This package provides an algorithmic environment fo describing algorithms.
% You can use the algorithmic environment in-text or within a figure
% environment to provide for a floating algorithm. Do NOT use the algorithm
% floating environment provided by algorithm.sty (by the same authors) or
% algorithm2e.sty (by Christophe Fiorio) as the IEEE does not use dedicated
% algorithm float types and packages that provide these will not provide
% correct IEEE style captions. The latest version and documentation of
% algorithmic.sty can be obtained at:
% http://www.ctan.org/pkg/algorithms
% Also of interest may be the (relatively newer and more customizable)
% algorithmicx.sty package by Szasz Janos:
% http://www.ctan.org/pkg/algorithmicx




% *** ALIGNMENT PACKAGES ***
%
%\usepackage{array}
% Frank Mittelbach's and David Carlisle's array.sty patches and improves
% the standard LaTeX2e array and tabular environments to provide better
% appearance and additional user controls. As the default LaTeX2e table
% generation code is lacking to the point of almost being broken with
% respect to the quality of the end results, all users are strongly
% advised to use an enhanced (at the very least that provided by array.sty)
% set of table tools. array.sty is already installed on most systems. The
% latest version and documentation can be obtained at:
% http://www.ctan.org/pkg/array


% IEEEtran contains the IEEEeqnarray family of commands that can be used to
% generate multiline equations as well as matrices, tables, etc., of high
% quality.




% *** SUBFIGURE PACKAGES ***
%\ifCLASSOPTIONcompsoc
%  \usepackage[caption=false,font=normalsize,labelfont=sf,textfont=sf]{subfig}
%\else
%  \usepackage[caption=false,font=footnotesize]{subfig}
%\fi
% subfig.sty, written by Steven Douglas Cochran, is the modern replacement
% for subfigure.sty, the latter of which is no longer maintained and is
% incompatible with some LaTeX packages including fixltx2e. However,
% subfig.sty requires and automatically loads Axel Sommerfeldt's caption.sty
% which will override IEEEtran.cls' handling of captions and this will result
% in non-IEEE style figure/table captions. To prevent this problem, be sure
% and invoke subfig.sty's "caption=false" package option (available since
% subfig.sty version 1.3, 2005/06/28) as this is will preserve IEEEtran.cls
% handling of captions.
% Note that the Computer Society format requires a larger sans serif font
% than the serif footnote size font used in traditional IEEE formatting
% and thus the need to invoke different subfig.sty package options depending
% on whether compsoc mode has been enabled.
%
% The latest version and documentation of subfig.sty can be obtained at:
% http://www.ctan.org/pkg/subfig




% *** FLOAT PACKAGES ***
%
%\usepackage{fixltx2e}
% fixltx2e, the successor to the earlier fix2col.sty, was written by
% Frank Mittelbach and David Carlisle. This package corrects a few problems
% in the LaTeX2e kernel, the most notable of which is that in current
% LaTeX2e releases, the ordering of single and double column floats is not
% guaranteed to be preserved. Thus, an unpatched LaTeX2e can allow a
% single column figure to be placed prior to an earlier double column
% figure.
% Be aware that LaTeX2e kernels dated 2015 and later have fixltx2e.sty's
% corrections already built into the system in which case a warning will
% be issued if an attempt is made to load fixltx2e.sty as it is no longer
% needed.
% The latest version and documentation can be found at:
% http://www.ctan.org/pkg/fixltx2e


%\usepackage{stfloats}
% stfloats.sty was written by Sigitas Tolusis. This package gives LaTeX2e
% the ability to do double column floats at the bottom of the page as well
% as the top. (e.g., "\begin{figure*}[!b]" is not normally possible in
% LaTeX2e). It also provides a command:
%\fnbelowfloat
% to enable the placement of footnotes below bottom floats (the standard
% LaTeX2e kernel puts them above bottom floats). This is an invasive package
% which rewrites many portions of the LaTeX2e float routines. It may not work
% with other packages that modify the LaTeX2e float routines. The latest
% version and documentation can be obtained at:
% http://www.ctan.org/pkg/stfloats
% Do not use the stfloats baselinefloat ability as the IEEE does not allow
% \baselineskip to stretch. Authors submitting work to the IEEE should note
% that the IEEE rarely uses double column equations and that authors should try
% to avoid such use. Do not be tempted to use the cuted.sty or midfloat.sty
% packages (also by Sigitas Tolusis) as the IEEE does not format its papers in
% such ways.
% Do not attempt to use stfloats with fixltx2e as they are incompatible.
% Instead, use Morten Hogholm'a dblfloatfix which combines the features
% of both fixltx2e and stfloats:
%
% \usepackage{dblfloatfix}
% The latest version can be found at:
% http://www.ctan.org/pkg/dblfloatfix




%\ifCLASSOPTIONcaptionsoff
%  \usepackage[nomarkers]{endfloat}
% \let\MYoriglatexcaption\caption
% \renewcommand{\caption}[2][\relax]{\MYoriglatexcaption[#2]{#2}}
%\fi
% endfloat.sty was written by James Darrell McCauley, Jeff Goldberg and 
% Axel Sommerfeldt. This package may be useful when used in conjunction with 
% IEEEtran.cls'  captionsoff option. Some IEEE journals/societies require that
% submissions have lists of figures/tables at the end of the paper and that
% figures/tables without any captions are placed on a page by themselves at
% the end of the document. If needed, the draftcls IEEEtran class option or
% \CLASSINPUTbaselinestretch interface can be used to increase the line
% spacing as well. Be sure and use the nomarkers option of endfloat to
% prevent endfloat from "marking" where the figures would have been placed
% in the text. The two hack lines of code above are a slight modification of
% that suggested by in the endfloat docs (section 8.4.1) to ensure that
% the full captions always appear in the list of figures/tables - even if
% the user used the short optional argument of \caption[]{}.
% IEEE papers do not typically make use of \caption[]'s optional argument,
% so this should not be an issue. A similar trick can be used to disable
% captions of packages such as subfig.sty that lack options to turn off
% the subcaptions:
% For subfig.sty:
% \let\MYorigsubfloat\subfloat
% \renewcommand{\subfloat}[2][\relax]{\MYorigsubfloat[]{#2}}
% However, the above trick will not work if both optional arguments of
% the \subfloat command are used. Furthermore, there needs to be a
% description of each subfigure *somewhere* and endfloat does not add
% subfigure captions to its list of figures. Thus, the best approach is to
% avoid the use of subfigure captions (many IEEE journals avoid them anyway)
% and instead reference/explain all the subfigures within the main caption.
% The latest version of endfloat.sty and its documentation can obtained at:
% http://www.ctan.org/pkg/endfloat
%
% The IEEEtran \ifCLASSOPTIONcaptionsoff conditional can also be used
% later in the document, say, to conditionally put the References on a 
% page by themselves.




% *** PDF, URL AND HYPERLINK PACKAGES ***
%
%\usepackage{url}
\usepackage{cite}
\usepackage{graphicx}
\usepackage{float}
% url.sty was written by Donald Arseneau. It provides better support for
% handling and breaking URLs. url.sty is already installed on most LaTeX
% systems. The latest version and documentation can be obtained at:
% http://www.ctan.org/pkg/url
% Basically, \url{my_url_here}.




% *** Do not adjust lengths that control margins, column widths, etc. ***
% *** Do not use packages that alter fonts (such as pslatex).         ***
% There should be no need to do such things with IEEEtran.cls V1.6 and later.
% (Unless specifically asked to do so by the journal or conference you plan
% to submit to, of course. )


% correct bad hyphenation here
\hyphenation{op-tical net-works semi-conduc-tor}


\begin{document}

%
% paper title
% Titles are generally capitalized except for words such as a, an, and, as,
% at, but, by, for, in, nor, of, on, or, the, to and up, which are usually
% not capitalized unless they are the first or last word of the title.
% Linebreaks \\ can be used within to get better formatting as desired.
% Do not put math or special symbols in the title.
\title{FireSmokeGenNet: Boundary-Aware Diffusion Synthesis and Multimodal Quality Filtering for Wildfire-Smoke Detection}
% Title moderated per reviewer request: the experiments in this manuscript
% evaluate still-image smoke detection/synthesis and do not measure
% detection delay, time-to-alarm, or other early-monitoring quantities.
% "Detection" is used instead of "Monitoring" until a dedicated
% early-detection/time-to-alarm experiment (Section VII, item F5) is added.
%
%
% author names and IEEE memberships
% note positions of commas and nonbreaking spaces ( ~ ) LaTeX will not break
% a structure at a ~ so this keeps an author's name from being broken across
% two lines.
% use \thanks{} to gain access to the first footnote area
% a separate \thanks must be used for each paragraph as LaTeX2e's \thanks
% was not built to handle multiple paragraphs
%

\author{Ammar Amjad,~\IEEEmembership{Member,~IEEE,}%
\thanks{Ammar Amjad is with the Department of Electrical and Computer Engineering, National Yang Ming Chiao Tung University, Hsinchu, Taiwan (e-mail: ammar@nycu.edu.tw).}%
% NOTE: "\and" is only valid in IEEEtran's conference/peerreviewca modes; in
% journal mode (used here) it silently fails to insert any separator or
% spacing, which previously rendered as "IEEE,Kenny Lin" with no space. A
% plain comma-separated list (the standard IEEEtran journal-mode author
% format) is used instead.
~Kenny Lin,~\IEEEmembership{Senior Engineer, Gama Innovation Inc.}%
\thanks{Kenny Lin is with the Department of Product Development and Engineering, Gama Innovation Inc., Taipei City, Taiwan (e-mail: kenneylin@hotmail.com).}}


% <-this % stops a space
%\thanks{M. Shell was with the Department
%of Electrical and Computer Engineering, Georgia Institute of Technology, Atlanta,
%GA, 30332 USA e-mail: (see http://www.michaelshell.org/contact.html).}% <-this % stops a space
%\thanks{J. Doe and J. Doe are with Anonymous University.}% <-this % stops a space
%\thanks{Manuscript received April 19, 2005; revised August 26, 2015.}}

% note the % following the last \IEEEmembership and also \thanks - 
% these prevent an unwanted space from occurring between the last author name
% and the end of the author line. i.e., if you had this:
% 
% \author{....lastname \thanks{...} \thanks{...} }
%                     ^------------^------------^----Do not want these spaces!
%
% a space would be appended to the last name and could cause every name on that
% line to be shifted left slightly. This is one of those "LaTeX things". For
% instance, "\textbf{A} \textbf{B}" will typeset as "A B" not "AB". To get
% "AB" then you have to do: "\textbf{A}\textbf{B}"
% \thanks is no different in this regard, so shield the last } of each \thanks
% that ends a line with a % and do not let a space in before the next \thanks.
% Spaces after \IEEEmembership other than the last one are OK (and needed) as
% you are supposed to have spaces between the names. For what it is worth,
% this is a minor point as most people would not even notice if the said evil
% space somehow managed to creep in.



% The paper headers
\markboth{IEEE Transactions on Multimedia}%
{Shell \MakeLowercase{\textit{et al.}}: Bare Demo of IEEEtran.cls for IEEE Journals}
% The only time the second header will appear is for the odd numbered pages
% after the title page when using the twoside option.
% 
% *** Note that you probably will NOT want to include the author's ***
% *** name in the headers of peer review papers.                   ***
% You can use \ifCLASSOPTIONpeerreview for conditional compilation here if
% you desire.




% If you want to put a publisher's ID mark on the page you can do it like
% this:
%\IEEEpubid{0000--0000/00\$00.00~\copyright~2015 IEEE}
% Remember, if you use this you must call \IEEEpubidadjcol in the second
% column for its text to clear the IEEEpubid mark.



% use for special paper notices
%\IEEEspecialpapernotice{(Invited Paper)}




% make the title area
\maketitle

% As a general rule, do not put math, special symbols or citations
% in the abstract or keywords.
\begin{abstract}
Smoke is typically the earliest visible indicator of a wildfire. With the rapid progress in deep learning, image-based smoke detection has become a vital approach for early wildfire recognition and prevention. However, a key bottleneck in training effective detection models is the limited availability of annotated forest fire smoke images. While generative models offer a promising solution for synthesizing such data, conventional inpainting approaches often struggle to maintain visual fidelity and contextual consistency—especially at the interface between synthesized smoke and complex natural backgrounds.

To address these challenges, we propose \textbf{FireSmokeGenNet}, a novel framework for generating realistic and diverse forest fire smoke images. Our pipeline begins by using a pre-trained segmentation model and a multimodal vision-language model to extract smoke masks and generate descriptive captions. We then introduce a dual-branch architecture that incorporates both the mask and its masked image to guide the generative process. A key contribution of our framework is the \textit{Mask Random Difference Loss (MRDL)}, which enforces local consistency by introducing stochastic mask perturbations that improve boundary smoothness and transparency during training.


\end{abstract}

% Note that keywords are not normally used for peerreview papers.
\begin{IEEEkeywords}
Diffusion models, environmental monitoring, forest fire detection, generative AI, image inpainting, joint cross-attention, large language models, mask conditioning, multimodal guidance, and synthetic smoke generation.
\end{IEEEkeywords}






% For peer review papers, you can put extra information on the cover
% page as needed:
% \ifCLASSOPTIONpeerreview
% \begin{center} \bfseries EDICS Category: 3-BBND \end{center}
% \fi
%
% For peerreview papers, this IEEEtran command inserts a page break and
% creates the second title. It will be ignored for other modes.
\IEEEpeerreviewmaketitle
% Required packages:
% \usepackage{amsmath}
% \usepackage{amssymb}
% \usepackage{booktabs}
% \usepackage{tabularx}
% \usepackage{array}
%
% Define the proposed method in the preamble:
\newcommand{\method}{FireSmokeGenNet}

\section{Introduction}
\label{sec:introduction}

\IEEEPARstart{W}{ildfires} cause extensive environmental, economic,
and societal damage, with severe fire seasons affecting millions of
hectares and releasing substantial quantities of carbon into the
atmosphere
\cite{Giglio2018-ca,Jain2024,Byrne2024}. Timely wildfire detection is
therefore important for early assessment, containment planning, and
emergency response. Smoke is frequently observable before flames become
clearly visible and is consequently an important visual indicator for
early wildfire detection
\cite{7823512,BOUGUETTAYA2022108309}. Unlike flames, which may be
partially occluded by terrain or vegetation, smoke can rise above the
forest canopy and remain visible over comparatively long distances.
Vision-based smoke detection has thus become an important component of
wildfire-monitoring systems, complementing satellite observations,
unmanned aerial vehicle surveillance, and ground-based sensor networks
\cite{BOUGUETTAYA2022108309,Karger2026}.

Deep object detectors, including YOLO, DETR, and EfficientDet, have
substantially advanced real-time object localization
\cite{redmon2016lookonceunifiedrealtime,
carion2020endtoendobjectdetectiontransformers,
tan2020efficientdetscalableefficientobject}. Their effectiveness,
however, depends strongly on the availability and diversity of
annotated training data. These requirements are difficult to satisfy
for early wildfire-smoke detection because real wildfire events are
geographically dispersed, temporally unpredictable, and hazardous to
observe at close range. Existing datasets provide valuable imagery for
detection and segmentation research but may represent a restricted
range of fire stages, camera viewpoints, illumination conditions,
seasons, and weather patterns
\cite{shamsoshoara2020aerialimagerypileburn}. Early-stage smoke is
especially difficult to collect and annotate because it may appear as
a faint, semi-transparent structure with low contrast against clouds,
haze, fog, vegetation, or bright sky. Models trained on restricted
environmental distributions may consequently exhibit reduced
robustness under previously unseen conditions, consistent with the
broader challenge of out-of-distribution generalization
\cite{hendrycks2021facesrobustnesscriticalanalysis}.

Synthetic data generation provides a scalable approach for increasing
the diversity of smoke appearances and environmental conditions
represented during detector training. Generating convincing smoke,
however, differs substantially from synthesizing rigid objects. Smoke
has no fixed geometry, contains structures at multiple spatial scales,
and exhibits spatially varying opacity. Its appearance depends on
illumination, background texture, atmospheric conditions, and plume
density. Moreover, its boundaries are naturally diffuse and irregular.
A useful synthetic-smoke generator must therefore preserve the
environmental context while producing diverse smoke structures,
plausible translucency, and smooth transitions around uncertain plume
boundaries.

Early synthetic-smoke methods relied on computational simulation,
procedural generation, and image compositing
\cite{6053841,6425498}. Such approaches offer explicit
control over smoke location and appearance but may require substantial
parameter adjustment and can retain a visual gap between simulated and
real smoke. GAN-based methods have also been investigated for
wildfire-data augmentation and controllable smoke generation
\cite{rs12223715,9249044}. Although GANs offer comparatively efficient
inference, their adversarial optimization may be affected by training
instability and insufficient mode coverage, particularly when modeling
fine turbulent structures and spatially varying opacity.

Diffusion models have recently achieved strong performance in
high-fidelity image synthesis. Denoising diffusion probabilistic models
learn to reconstruct data by reversing a progressive Gaussian
corruption process
\cite{ho2020denoisingdiffusionprobabilisticmodels}, while improved
formulations introduce more effective noise schedules and
reverse-process variance parameterizations
\cite{nichol2021improveddenoisingdiffusionprobabilistic}. Latent
diffusion models reduce computational requirements by performing the
denoising process in a compressed representation space and incorporate
external conditions using cross-attention
\cite{rombach2022highresolutionimagesynthesislatent}. Spatial-control
architectures such as ControlNet and T2I-Adapter further demonstrate
that external structural signals can be injected into multiple stages
of a pretrained diffusion model
\cite{zhang2023addingconditionalcontroltexttoimage,
mou2023t2iadapterlearningadaptersdig}. For wildfire imagery,
FlameDiffuser employs mask-guided latent diffusion to generate synthetic
fire scenes for data augmentation
\cite{wang2024flamediffuserwildfireimage}.

Despite this progress, three limitations remain important for
diffusion-based wildfire-smoke synthesis. First, standard inpainting
commonly combines the noisy latent, binary mask, and masked-image
representation through input-level concatenation. This formulation does
not explicitly separate structural information from contextual
appearance. Second, deterministic binary masks provide rigid spatial
constraints that do not represent the uncertain and semi-transparent
transitions observed at smoke boundaries. Third, diffusion sampling
produces samples with varying perceptual and task-level quality, and
indiscriminately adding all generated images to detector training can
introduce unrealistic colors, insufficient visibility, or implausible
opacity. Existing wildfire-synthesis pipelines have not jointly
addressed these limitations through decoupled multiscale conditioning,
boundary-specific consistency regularization, and task-aware synthetic
sample filtering.

To address these limitations, we propose \method, a controllable
latent-diffusion framework for synthetic wildfire-smoke generation.
The framework separately encodes the target smoke mask and masked-image
context, injects their multiscale representations hierarchically into
the diffusion U-Net, and combines them through Joint Cross-Attention
(JCA). We further introduce the Mask Random Difference Loss (MRDL), a
stochastic consistency objective that regularizes the model around
morphologically perturbed smoke boundaries. Finally, a task-specific
vision--language quality-assessment module ranks generated samples
according to smoke color, visibility, and translucency before their use
in downstream detector training.

The main contributions of this work are summarized as follows:

\begin{itemize}
    \item We propose a hierarchical dual-branch conditioning
    architecture that represents smoke-mask geometry and masked-image
    context through separate feature streams. The proposed JCA module
    independently attends to the structural and contextual features
    before fusing their representations at intermediate U-Net
    resolutions.

    \item We introduce MRDL, a stochastic boundary-consistency
    regularizer inspired by prediction consistency under
    semantics-preserving perturbations
    \cite{sajjadi2016regularizationstochastictransformationsperturbations}.
    MRDL penalizes differences between denoising predictions conditioned
    on the original and morphologically perturbed masks within the
    affected latent-space boundary region.

    \item We develop a task-specific multimodal quality-filtering
    strategy based on Qwen2-VL-7B
    \cite{wang2024qwen2vlenhancingvisionlanguagemodels}. Motivated by
    learned human-preference assessment
    \cite{xu2023imagerewardlearningevaluatinghuman,
    wu2023humanpreferencescorebetter}, the module evaluates generated
    smoke according to color fidelity, visibility, and translucency
    realism and retains the highest-ranked samples for downstream
    detector training.

    \item We establish an equal-budget evaluation protocol comparing
    real-only training, additional real data, randomly selected
    synthetic data, and VLM-filtered synthetic data. This protocol
    separates the benefit of synthetic-data generation from the
    contribution of quality filtering and evaluates robustness under
    environmental distribution shifts.
\end{itemize}

The remainder of this paper is organized as follows. Section~II reviews
wildfire-smoke datasets, synthetic smoke generation, diffusion-based
inpainting, and learned image-quality assessment. Section~III introduces
the required diffusion-model preliminaries and formulates the
smoke-boundary uncertainty problem. Section~IV presents the proposed
architecture, JCA module, MRDL objective, and multimodal filtering
pipeline. Section~V describes the datasets, implementation details,
evaluation protocols, and metrics. Section~VI reports the quantitative,
qualitative, domain-shift, and ablation results. Section~VII discusses
limitations and future directions, and Section~VIII concludes the
paper.
\section{Related Work}
\label{sec:related}


\subsection{Wildfire-Smoke Datasets and Data Augmentation}
\label{subsec:related_datasets}

The development of reliable wildfire-smoke detectors is closely tied
to the availability and diversity of annotated data. The FLAME dataset
provides aerial visible and thermal imagery acquired during prescribed
pile burns and supports fire classification and segmentation research
\cite{shamsoshoara2020aerialimagerypileburn}. Although FLAME is an
important benchmark for aerial fire analysis, its acquisition under a
limited set of controlled burning conditions restricts its
representation of environmental and geographical variability.

For early smoke detection from fixed-view cameras, Dewangan
\emph{et al.} introduced FIgLib, a collection of approximately 25,000
labeled images derived from wildfire-camera sequences in Southern
California, together with the spatiotemporal SmokeyNet architecture
\cite{Dewangan_2022}. The dataset includes smoke and non-smoke
observations captured before and after reported fire ignition.
Nevertheless, its geographical concentration and reliance on fixed
camera infrastructure may limit direct generalization to other regions,
camera configurations, and atmospheric conditions.

Several datasets have also been developed for smoke localization and
segmentation. Smoke100K contains synthetic smoke images paired with
smoke-free backgrounds, smoke masks, and bounding-box annotations
\cite{cheng_gcce19_smoke100k}. SMOKE5K combines real and synthetic images
with pixel-level smoke annotations and explicitly considers the
uncertainty associated with transparent and irregular smoke regions
\cite{yan2023transmissionguidedbayesiangenerativemodel}. These resources are valuable for detection
and segmentation; however, synthetic subsets constructed through
compositing can exhibit appearance differences from real smoke,
particularly around semi-transparent boundaries.

Data augmentation has consequently become an important strategy for
improving smoke-detector robustness. Conventional transformations,
including geometric, photometric, and color-space augmentation, improve
sample diversity but cannot introduce genuinely new smoke structures.
Compositing-based approaches offer greater control by blending
foreground smoke masks with smoke-free scenes
\cite{6425498}. Their realism,
however, depends strongly on the quality of the source masks, opacity
estimation, color adjustment, and boundary blending. These limitations
motivate generative approaches that learn smoke appearance directly
from data.


\subsection{Generative Models for Smoke and Wildfire Synthesis}
\label{subsec:related_generative_smoke}

GAN-based generation has been investigated as an alternative to
procedural simulation and deterministic compositing. Park
\emph{et al.} employed CycleGAN-based augmentation to improve wildfire
classification from remote-camera imagery
\cite{rs12223715}. Xie and Tao proposed controllable smoke-generation
networks that decompose an image into smoke and non-smoke components,
allowing smoke characteristics to be modified through latent
representations
\cite{9249044}. These studies demonstrated the potential of learned
generation for expanding smoke-related datasets.

Despite their efficient inference, GAN-based methods can be sensitive
to adversarial-training instability and insufficient mode coverage.
For smoke synthesis, these issues may appear as repetitive plume
structures, inconsistent opacity, or inadequate integration with the
background. Moreover, methods designed for image-level augmentation do
not necessarily preserve the spatial correspondence required for
automatically transferring masks or bounding-box annotations to
generated samples.

Diffusion models provide an alternative generative formulation based
on iterative denoising. DDPMs learn a reverse process that reconstructs
samples from progressively corrupted observations
\cite{ho2020denoisingdiffusionprobabilisticmodels}. Subsequent
developments improved noise scheduling, variance parameterization, and
sampling efficiency
\cite{nichol2021improveddenoisingdiffusionprobabilistic}. Latent
diffusion models perform the denoising process within a compressed
representation space and incorporate semantic conditions through
cross-attention
\cite{rombach2022highresolutionimagesynthesislatent}. Compared with
adversarial generators, diffusion models generally provide stable
training and strong sample diversity, although their iterative
inference incurs greater computational cost.


\subsection{Spatially Controlled Diffusion and Image Inpainting}
\label{subsec:related_controlled_diffusion}

Diffusion-based image inpainting uses masks, image context, text, or
other structural signals to constrain the denoising process. RePaint
conditions a pretrained unconditional diffusion model by repeatedly
resampling the known image regions during reverse diffusion, enabling
inpainting across different mask configurations without mask-specific
model retraining
\cite{lugmayr2022repaintinpaintingusingdenoising}. SmartBrush incorporates text and object-shape
guidance to improve spatial controllability and background preservation
during object inpainting
\cite{xie2022smartbrushtextshapeguided}. InstructPix2Pix learns instruction-guided
image editing from synthetically constructed editing pairs
\cite{brooks2023instructpix2pixlearningfollowimage}.

More general spatial-control architectures inject external conditions
into pretrained latent diffusion models. ControlNet introduces
trainable control branches connected to a frozen diffusion backbone,
allowing conditions such as edges, depth maps, poses, and segmentation
maps to guide generation
\cite{zhang2023addingconditionalcontroltexttoimage}. T2I-Adapter instead
uses lightweight adapters to align external control signals with the
internal representations of a pretrained text-to-image model
\cite{mou2023t2iadapterlearningadaptersdig}. These methods demonstrate
the value of hierarchical structural guidance but are designed for
general image synthesis rather than the uncertain, semi-transparent
boundaries characteristic of smoke.

For wildfire synthesis, FLAME Diffuser introduced a training-free,
mask-guided diffusion pipeline that uses wildfire masks and
Perlin-noise-based guidance to control the location and extent of
generated fire regions
\cite{wang2024flamediffuserwildfireimage}. It demonstrates that
diffusion-generated wildfire imagery can support downstream
data-augmentation tasks. However, its primary emphasis is flame
generation and paired spatial annotation rather than explicit modeling
of semi-transparent smoke interfaces.

In contrast, the present work focuses specifically on smoke synthesis.
It separates mask geometry from masked-background context through two
conditioning branches, fuses the resulting representations using Joint
Cross-Attention, and introduces a boundary-restricted consistency loss
for stochastic mask perturbations. This formulation is intended to
address smoke-boundary uncertainty while retaining spatial control and
environmental context.


\subsection{Learned Assessment of Generated-Image Quality}
\label{subsec:related_quality_assessment}

Generated-image evaluation has traditionally relied on distributional
and perceptual metrics such as Fréchet Inception Distance, structural
similarity, learned perceptual similarity, and text--image embedding
similarity. Although these measures capture complementary aspects of
image quality, they may not reflect task-specific requirements or
human judgments of realism.

Recent preference-based models learn quality assessments directly from
human comparisons. ImageReward uses expert-ranked text--image pairs to
learn a reward model for evaluating and improving text-to-image
generation
\cite{xu2023imagerewardlearningevaluatinghuman}. Human Preference Score
similarly adapts vision--language representations to predict human
preferences among generated images
\cite{wu2023humanpreferencescorebetter}. These approaches demonstrate
that learned multimodal scoring can complement conventional automatic
metrics.

Large vision--language models provide a further mechanism for
evaluating images according to explicitly defined semantic criteria.
Qwen2-VL supports joint visual and textual reasoning across variable
image resolutions
\cite{wang2024qwen2vlenhancingvisionlanguagemodels}. However,
general-purpose preference models are not specifically trained to
assess the color fidelity, visibility, and translucency of wildfire
smoke. Moreover, a generated image that appears perceptually appealing
is not necessarily useful for downstream smoke detection.

The proposed framework therefore adapts a vision--language model using
smoke-specific human annotations and evaluates the effect of its
filtering decisions through an equal-budget detector-training protocol.
Unlike general-purpose generated-image ranking, the proposed filtering
stage is explicitly evaluated according to its contribution to
downstream wildfire-smoke detection.
\section{Preliminaries}
\label{sec:prelim}


\subsection{Latent Diffusion Models}
\label{subsec:prelim_ldm}

Latent diffusion models (LDMs) perform the diffusion process in a
compressed representation space learned by a variational autoencoder
(VAE), thereby reducing the computational cost of high-resolution image
generation
\cite{rombach2022highresolutionimagesynthesislatent}. Given an RGB image

\begin{equation}
x\in\mathbb{R}^{H\times W\times3},
\end{equation}

the VAE encoder \(\mathcal{E}\) maps \(x\) to a latent representation

\begin{equation}
z_0=\mathcal{E}(x)
\in\mathbb{R}^{h\times w\times c},
\label{eq:vae_encoding}
\end{equation}

where \(h=H/f\), \(w=W/f\), and \(f\) is the spatial compression
factor. For the VAE adopted in this work, \(f=8\) and \(c=4\). After
reverse diffusion, the VAE decoder \(\mathcal{D}\) reconstructs the
output image from the denoised latent:

\begin{equation}
\widehat{x}=\mathcal{D}(\widehat{z}_0).
\label{eq:vae_decoding}
\end{equation}

The VAE parameters remain fixed during training of the conditional
diffusion model.


\subsection{Forward and Reverse Diffusion Processes}
\label{subsec:prelim_diffusion}

The diffusion formulation comprises a fixed forward corruption process
and a learned reverse denoising process. As illustrated in
Fig.~\ref{fig:diffusion}, the forward process progressively perturbs the
clean latent \(z_0\) with Gaussian noise. The reverse process predicts
and removes the injected noise under the supplied conditioning
information.

\begin{figure*}[t]
\centering
\includegraphics[width=\textwidth]{fig44.png}
\caption{Forward and reverse latent-diffusion processes. The forward
process progressively perturbs the clean latent representation \(z_0\)
with Gaussian noise. During reverse diffusion, the conditional U-Net
iteratively predicts and removes noise using the noisy latent, diffusion
timestep, and multimodal conditioning information. The final denoised
latent is decoded by the fixed VAE decoder to obtain the generated
image.}
\label{fig:diffusion}
\end{figure*}


\subsubsection{Forward Diffusion Process}

Let \(\{\beta_t\}_{t=1}^{T}\) denote a predefined variance schedule,
where \(0<\beta_t<1\). We define

\begin{equation}
\alpha_t=1-\beta_t,
\qquad
\bar{\alpha}_t=\prod_{\tau=1}^{t}\alpha_\tau.
\label{eq:alpha_definition}
\end{equation}

The forward Markov transition is

\begin{equation}
q(z_t\mid z_{t-1})
=
\mathcal{N}
\left(
z_t;
\sqrt{\alpha_t}\,z_{t-1},
\beta_t I
\right).
\label{eq:forward_transition}
\end{equation}

Using the reparameterization property of Gaussian distributions,
\(z_t\) can be sampled directly from \(z_0\):

\begin{equation}
z_t
=
\sqrt{\bar{\alpha}_t}\,z_0
+
\sqrt{1-\bar{\alpha}_t}\,\epsilon,
\qquad
\epsilon\sim\mathcal{N}(0,I).
\label{eq:forward}
\end{equation}

We employ the cosine noise schedule introduced by Nichol and Dhariwal
\cite{nichol2021improveddenoisingdiffusionprobabilistic}. Let

\begin{equation}
g(t)
=
\cos^2
\left(
\frac{t/T+s_0}{1+s_0}\,
\frac{\pi}{2}
\right),
\qquad s_0=0.008.
\label{eq:cosine_function}
\end{equation}

The cumulative noise schedule is normalized as

\begin{equation}
\bar{\alpha}_t
=
\frac{g(t)}{g(0)}.
\label{eq:cosine_schedule}
\end{equation}

The corresponding variance coefficients are obtained from successive
values of \(\bar{\alpha}_t\):

\begin{equation}
\beta_t
=
\min
\left(
1-
\frac{\bar{\alpha}_t}{\bar{\alpha}_{t-1}},
\beta_{\max}
\right),
\label{eq:cosine_beta}
\end{equation}

where \(\beta_{\max}<1\) prevents numerical instability near the final
diffusion steps.


\subsubsection{Conditional Reverse Diffusion Process}

The reverse process reconstructs the clean latent by iteratively
sampling

\begin{equation}
p_\theta(z_{t-1}\mid z_t,C)
=
\mathcal{N}
\left(
z_{t-1};
\mu_\theta(z_t,t,C),
\sigma_t^2 I
\right),
\label{eq:reverse}
\end{equation}

where \(C\) denotes the conditioning information and \(\theta\)
represents the trainable parameters of the diffusion U-Net. The
diffusion timestep \(t\) is provided separately through a timestep
embedding and is therefore not included in \(C\).

Under the noise-prediction parameterization, the reverse-process mean
is expressed as

\begin{equation}
\mu_\theta(z_t,t,C)
=
\frac{1}{\sqrt{\alpha_t}}
\left(
z_t-
\frac{1-\alpha_t}{\sqrt{1-\bar{\alpha}_t}}
\epsilon_\theta(z_t,t,C)
\right).
\label{eq:reverse_mean}
\end{equation}

where \(\epsilon_\theta(z_t,t,C)\) is the noise predicted by the
conditional U-Net. Depending on the adopted DDPM parameterization, the
reverse variance \(\sigma_t^2\) may be fixed according to the posterior
variance or learned during training
\cite{nichol2021improveddenoisingdiffusionprobabilistic}. The specific
parameterization used in FireSmokeGenNet is reported in
Section~\ref{sec:method}.


\subsubsection{Classifier-Free Guidance}

Classifier-free guidance combines conditional and unconditional noise
predictions without requiring an external classifier
\cite{ho2022classifierfreediffusionguidance}. During training, the conditioning information
is randomly replaced by a null condition. At inference time, the guided
noise prediction is

\begin{equation}
\begin{aligned}
\widetilde{\epsilon}_{\theta}(z_t,t,C)
={}&
\epsilon_{\theta}(z_t,t,\varnothing) \\
&+
\gamma
\left[
\epsilon_{\theta}(z_t,t,C)
-
\epsilon_{\theta}(z_t,t,\varnothing)
\right],
\end{aligned}
\label{eq:classifier_free_guidance}
\end{equation}

where \(\gamma\geq1\) is the guidance scale and \(\varnothing\)
denotes the null condition. Increasing \(\gamma\) generally strengthens
adherence to the conditioning information but may reduce sample
diversity or introduce visual artifacts.


\subsubsection{Diffusion Training Objective}

The standard noise-prediction objective is

\begin{equation}
\begin{aligned}
\mathcal{L}_{\mathrm{diff}}
=
\mathbb{E}_{
\substack{
x\sim p_{\mathrm{data}},\,
t\sim\mathcal{U}\{1,\ldots,T\},\\
\epsilon\sim\mathcal{N}(0,I)
}
}
\left[
\left\|
\epsilon-
\epsilon_\theta(z_t,t,C)
\right\|_2^2
\right],
\end{aligned}
\label{eq:loss_diff}
\end{equation}

where \(z_0=\mathcal{E}(x)\) and \(z_t\) is generated using
Eq.~\eqref{eq:forward}. This simplified objective corresponds to a
reweighted form of the variational objective used to optimize the
reverse diffusion process
\cite{ho2020denoisingdiffusionprobabilisticmodels}.


\subsubsection{Conditional Score Interpretation}

The noise-prediction network is related to the conditional score of the
perturbed latent distribution:

\begin{equation}
\nabla_{z_t}
\log p(z_t\mid C)
\approx
-
\frac{
\epsilon_\theta(z_t,t,C)
}{
\sqrt{1-\bar{\alpha}_t}
}.
\label{eq:conditional_score}
\end{equation}

This interpretation connects denoising diffusion models with
score-based generative modeling
\cite{song2021scorebasedgenerativemodelingstochastic}. The conditioning information therefore guides the
estimated score toward latent configurations that are compatible with
the specified generation constraints.


\subsection{Conventional Latent-Inpainting Conditioning}
\label{subsec:flat_conditioning}

Let \(M\in\{0,1\}^{H\times W}\) denote a binary mask, where
\(M_{ij}=1\) identifies the region to be synthesized. The corresponding
masked image is

\begin{equation}
x_m=x\odot(1-M),
\label{eq:masked_image}
\end{equation}

where \(\odot\) denotes elementwise multiplication. The masked-image
latent is obtained as

\begin{equation}
z_m=\mathcal{E}(x_m)
\in\mathbb{R}^{h\times w\times c}.
\label{eq:masked_latent}
\end{equation}

The mask is resized to the latent spatial resolution:

\begin{equation}
M_z
=
\operatorname{Resize}(M,h,w)
\in\{0,1\}^{h\times w\times1}.
\label{eq:latent_mask}
\end{equation}

A conventional latent-inpainting U-Net forms its input through
channelwise concatenation:

\begin{equation}
u_t
=
[z_t;z_m;M_z]
\in
\mathbb{R}^{h\times w\times(2c+1)}.
\label{eq:flat_cond}
\end{equation}

For \(c=4\), this produces a nine-channel input. This formulation is
computationally convenient and provides direct spatial conditioning.
However, after input-level concatenation, it does not explicitly retain
separate structural and contextual conditioning pathways.


\subsection{Limitations of Flat Conditioning for Smoke Synthesis}
\label{subsec:flat_conditioning_limitations}

Although input-level concatenation is effective for general image
inpainting, smoke synthesis introduces several additional requirements.

\subsubsection{Implicit Multiscale Conditioning}

The binary mask and masked-image latent are directly available at the
input layer, whereas deeper U-Net blocks receive them through
progressively transformed intermediate representations. The
architecture therefore does not provide explicit multiscale access to
the original structural and contextual features. This may make it more
difficult to preserve global smoke geometry while refining local
texture, opacity, and boundary transitions.

\subsubsection{Heterogeneous Conditioning Sources}

The noisy latent \(z_t\), masked-image latent \(z_m\), and binary mask
\(M_z\) represent different types of information. The first two are
continuous latent representations, whereas the mask is a discrete
spatial constraint. Direct concatenation requires the initial U-Net
layers to learn how to separate and integrate these heterogeneous
signals without modality-specific encoders.

This does not imply that concatenation is intrinsically ineffective.
Rather, it motivates investigating whether separate structural and
contextual encoders can provide more explicit conditioning for smoke
generation.

\subsubsection{Deterministic Boundary Representation}

A binary mask defines a sharp distinction between the synthesis and
preservation regions. Real smoke boundaries, however, are often
semi-transparent, irregular, and difficult to localize precisely.
Treating a single binary contour as an exact boundary may therefore
introduce an overly rigid spatial assumption. Depending on the learned
behavior, this can produce abrupt transitions near the mask boundary or
excessive smoke leakage outside the intended region.

\subsubsection{Static Input-Level Representation}

The mask and masked-image inputs remain fixed throughout the reverse
process. Although timestep embeddings allow the U-Net to learn
noise-level-dependent behavior, conventional concatenation does not
explicitly control how structural and contextual information is
distributed across different U-Net resolutions. This observation
motivates the hierarchical conditioning strategy introduced in
Section~\ref{sec:method}.


\subsection{Stochastic Smoke-Boundary Formulation}
\label{subsec:boundary_uncertainty}

For smoke synthesis, the binary mask \(M\) is treated as an approximate
spatial constraint rather than an exact physical boundary. We define a
conditional perturbation distribution

\begin{equation}
M'\sim p_{\mathcal{T}}(M'\mid M),
\label{eq:mask_perturbation_distribution}
\end{equation}

where \(p_{\mathcal{T}}(M'\mid M)\) represents local morphological
variations of \(M\), such as controlled dilation and erosion. These
perturbations do not simulate the physical evolution of smoke. Instead,
they represent uncertainty in the position of the specified or
annotated smoke boundary.

Let \(C(M)\) and \(C(M')\) denote conditioning representations obtained
from the original and perturbed masks, respectively. For the same noisy
latent \(z_t\), noise realization, and timestep \(t\), the sensitivity
of the noise predictor to the mask perturbation can be expressed as

\begin{equation}
\Delta_\theta(M,M')
=
\epsilon_\theta(z_t,t,C(M))
-
\epsilon_\theta(z_t,t,C(M')).
\label{eq:boundary_prediction_difference}
\end{equation}

A general boundary-consistency objective can then be written as

\begin{equation}
\begin{aligned}
\mathcal{R}_{\mathrm{boundary}}
=
\mathbb{E}_{M'\sim p_{\mathcal{T}}(M'\mid M)}
\Big[
d_B\big(
&\epsilon_\theta(z_t,t,C(M)),\\
&\epsilon_\theta(z_t,t,C(M'))
\big)
\Big],
\end{aligned}
\label{eq:boundary_regularization_general}
\end{equation}

where \(d_B(\cdot,\cdot)\) measures prediction differences within the
boundary region affected by the perturbation. This expectation can be
estimated through stochastic mask sampling during training.

Equation~\eqref{eq:boundary_regularization_general} defines a
consistency objective over plausible mask perturbations; it is not a
direct likelihood marginalization over physical smoke boundaries. The
proposed MRDL provides a boundary-restricted and normalized
instantiation of \(d_B(\cdot,\cdot)\), as described in
Section~\ref{sec:method}.
\section{Methodology}
\label{sec:method}

\subsection{Overall Architecture}

\method~extends the SD-2 Inpainting U-Net \cite{rombach2022highresolutionimagesynthesislatent} with three major modifications: (1) a dual-branch conditioning encoder, (2) joint cross-attention (JCA) for feature fusion, and (3) stochastic mask perturbation via MRDL. Figure~\ref{fig:architecture} illustrates the complete pipeline.

\begin{figure*}[t]
\centering
\includegraphics[width=\textwidth]{fig11.png}
\caption{\textbf{Overall architecture of FireSmokeGenNet.} The proposed framework integrates multimodal guidance, dual-branch feature encoding, Joint Cross-Attention (JCA), conditional diffusion generation, and the proposed Mask Random Difference Loss (MRDL) to synthesize realistic forest fire smoke images with improved boundary consistency and contextual fidelity.}
\label{fig:architecture}
\end{figure*}

\subsection{Dual-Branch Conditioning Encoder}

To effectively capture complementary structural and contextual information, the proposed framework employs a dual-branch conditioning encoder, as illustrated in Fig.~\ref{fig:dual_encoder}. The binary smoke mask and masked background image are processed independently using dedicated backbone networks. A lightweight ResNet-18 encoder is adopted for the mask branch to efficiently extract structural and boundary information from the binary smoke mask, whereas a ResNet-50 encoder is employed for the masked RGB image to capture rich semantic and contextual representations. The extracted features are projected into a common latent space and concatenated before being forwarded to the proposed Joint Cross-Attention (JCA) module for multimodal feature fusion.

\begin{figure*}[t]
\centering
\includegraphics[width=\textwidth]{fig22.png}
\caption{\textbf{Dual-Branch Conditioning Encoder.} The proposed dual-branch conditioning encoder separately extracts structural and semantic information from the smoke mask and masked image. A lightweight ResNet-18 backbone encodes the binary smoke mask, while a ResNet-50 backbone extracts contextual features from the masked RGB image. Multi-scale feature representations are projected into a common latent space and concatenated to provide complementary conditioning information for the subsequent Joint Cross-Attention (JCA) module.}
\label{fig:dual_encoder}
\end{figure*}

\paragraph{Feature Extraction.}
Given a binary smoke mask $M$ and its corresponding masked image $x_m$, the mask branch extracts structural representations using a frozen ResNet-18 backbone, while the image branch employs a frozen ResNet-50 backbone to preserve rich semantic information learned from ImageNet pretraining. Multi-scale feature maps are extracted from both branches to facilitate hierarchical conditioning of the diffusion U-Net.

For the image branch, the extracted feature maps are

\begin{align}
F_I^{(1)} &\in \mathbb{R}^{H/2 \times W/2 \times 64},\\
F_I^{(2)} &\in \mathbb{R}^{H/4 \times W/4 \times 256},\\
F_I^{(3)} &\in \mathbb{R}^{H/8 \times W/8 \times 512},\\
F_I^{(4)} &\in \mathbb{R}^{H/16 \times W/16 \times 1024},\\
F_I^{(5)} &\in \mathbb{R}^{H/32 \times W/32 \times 2048},
\end{align}

whereas the mask branch generates corresponding multi-scale representations

\begin{equation}
F_M^{(s)}, \quad s \in \{2,3,4,5\},
\end{equation}

which capture the geometric structure of the smoke regions.

\paragraph{Feature Projection.}
Since the extracted feature dimensions differ from those required by the diffusion U-Net, each feature map is projected into the corresponding latent dimension using lightweight learnable $1\times1$ convolutional layers,

\begin{equation}
\hat{F}^{(s)}=\mathrm{Proj}^{(s)}\!\left(F^{(s)}\right),
\end{equation}

where

\[
\hat{F}^{(s)}\in
\mathbb{R}^{B\times C_{\mathrm{U\text{-}Net}}^{(s)}\times H_s\times W_s},
\]

and $C_{\mathrm{U\text{-}Net}}^{(s)}\in\{320,640,1280,1280\}$ denotes the channel dimension of the corresponding U-Net stage. These projected feature maps provide complementary structural and semantic guidance for the subsequent Joint Cross-Attention module while keeping the backbone encoders frozen throughout training.
\subsection{Spatially Adaptive Feature Injection}

Crucially, we inject features \emph{hierarchically} based on semantic granularity. Table~\ref{tab:injection} specifies the injection strategy for each U-Net resolution block.

\begin{table}[t]
\centering
\caption{\textbf{Spatially adaptive feature injection strategy.}}
\label{tab:injection}
\begin{tabular}{lcccc}
\toprule
U-Net Block & Resolution & Feature Source & Injection Type \\
\midrule
Down block 1 & $H/32 \times W/32$ & $F_m$ only & Cross-attention \\
Down block 2 & $H/16 \times W/16$ & $F_m$ only & Cross-attention \\
Down block 3 & $H/8 \times W/8$ & $F_m + F_M$ & JCA \\
Bottleneck & $H/8 \times W/8$ & $F_m$ only & Cross-attention \\
Up block 1 & $H/8 \times W/8$ & $F_m + F_M$ & JCA \\
Up block 2 & $H/16 \times W/16$ & $F_M$ only & Cross-attention \\
Up block 3 & $H/32 \times W/32$ & $F_M$ only & Cross-attention \\
Up block 4 & $H/64 \times W/64$ & $F_M$ only & Cross-attention \\
\bottomrule
\end{tabular}
\end{table}

\paragraph{Justification.}
This design respects the information hierarchy:
\begin{itemize}
    \item \textbf{Low resolutions ($H/32$, $H/16$):} At these coarse scales, only mask features are needed to specify where smoke should appear. Image features would add distracting texture details at a stage where structure is being formed.
    \item \textbf{Mid resolutions ($H/8$):} Both mask and image features are necessary. The mask provides structural constraints, while the image provides context for blending (e.g., matching background colors).
    \item \textbf{High resolutions ($H/16$, $H/32$, $H/64$):} Only image features are injected. At these fine scales, the mask is already satisfied; the model focuses on texture synthesis and boundary refinement.
\end{itemize}

\subsection{Joint Cross-Attention (JCA)}

To effectively integrate structural information from the smoke mask with semantic context extracted from the masked image, we propose the Joint Cross-Attention (JCA) module, as illustrated in Fig.~\ref{fig:jca}. Unlike conventional cross-attention that relies on a single conditioning source, JCA independently models the interactions between the U-Net latent features and the mask and image representations. The resulting attention features are subsequently fused to provide richer multimodal guidance for smoke synthesis.

\begin{figure*}[t]
\centering
\includegraphics[width=\textwidth]{fig33.png}
\caption{\textbf{Architecture of the proposed Joint Cross-Attention (JCA) module.} Multi-scale mask features $F_M$ and image features $F_I$ are independently projected into query, key, and value representations using learnable linear layers. Scaled dot-product attention is performed to capture complementary structural and semantic information from both modalities. The attended features are fused to generate the joint representation $F_J$, which is injected into the diffusion U-Net to guide realistic smoke generation.}
\label{fig:jca}
\end{figure*}

For layers requiring both modalities, we propose Joint Cross-Attention (JCA). Standard cross-attention computes

\begin{equation}
\text{CrossAttn}(Q,K,V)=
\operatorname{Softmax}\left(\frac{QK^{T}}{\sqrt{d}}\right)V
\label{eq:cross_attn}
\end{equation}

where $Q$ is derived from the U-Net latent features, while $K$ and $V$ are obtained from the conditioning features.

JCA computes separate attention for each conditioning branch and fuses the resulting representations through a lightweight multilayer perceptron (MLP):

\begin{align}
Q &= \text{Linear}_{Q}(X)\in\mathbb{R}^{B\times N\times d},\\
K_m &= \text{Linear}_{K}(\hat{F}_m), \qquad
V_m = \text{Linear}_{V}(\hat{F}_m),\\
K_M &= \text{Linear}_{K}(\hat{F}_M), \qquad
V_M = \text{Linear}_{V}(\hat{F}_M),\\
Z_m &= \text{CrossAttn}(Q,K_m,V_m),\\
Z_M &= \text{CrossAttn}(Q,K_M,V_M),\\
X' &= X+\text{MLP}([Z_m;Z_M]),
\label{eq:jca}
\end{align}

where $[\cdot;\cdot]$ denotes channel-wise concatenation, and the MLP consists of two fully connected layers with a hidden dimension of $4d$.

\paragraph{Computational Complexity}
Standard cross-attention has computational complexity $\mathcal{O}(N^{2}+NL)$, where $N$ denotes the number of spatial tokens and $L$ represents the conditioning sequence length. Since JCA computes attention for both the mask and image branches, its complexity becomes $\mathcal{O}(N^{2}+2NL)$. In practice, the additional computational overhead is modest because the conditioning sequence is relatively short (e.g., $L\approx1024$ for a $32\times32$ feature map), while the improved multimodal feature fusion significantly enhances the quality and structural consistency of the generated smoke.
\subsection{Mask Random Difference Loss (MRDL)}

MRDL addresses the deterministic boundary assumption by training the model to be invariant to plausible mask perturbations.

\paragraph{Morphological Perturbation Generation.}

Figure~\ref{fig:mrdl_mechanism} illustrates the stochastic
perturbation process used to represent uncertainty around the
annotated smoke boundary.

\begin{figure}[t]
\centering
\includegraphics[width=\linewidth]{fig21.png}
\caption{Mask Random Difference Loss (MRDL) mechanism.
Random morphological perturbations represent uncertainty in the
annotated smoke boundary and encourage consistent predictions across
plausible mask variations.}
\label{fig:mrdl_mechanism}
\end{figure}

For each training iteration, a perturbed mask $M'$ is generated by
applying a sequence of random morphological operations:

\begin{equation}
M'
=
\operatorname{Morph}(M,k,N),
\label{eq:mask_perturbation}
\end{equation}
where $k$ and $N$ are drawn independently from \emph{discrete} uniform
distributions, $k\sim\mathcal{U}\{10,11,\ldots,20\}$ (pixels) and
$N\sim\mathcal{U}\{1,2,3\}$ (operations), and $k$ denotes the
\emph{kernel radius}: with the neighborhood $u,v\in[-k,k]$ used in
Eqs.~\eqref{eq:dilation}--\eqref{eq:erosion} below, the resulting square
structuring element spans $(2k+1)\times(2k+1)$ pixels (e.g., $k=10$
corresponds to a $21\times21$ structuring element, and $k=20$ to a
$41\times41$ structuring element). At each perturbation step, dilation or
erosion is selected with equal probability.

Using a square structuring element, dilation is defined as

\begin{equation}
\begin{aligned}
\bigl[
\operatorname{Dilate}(M,k)
\bigr]_{ij}
=
\max_{\substack{-k\leq u\leq k\\-k\leq v\leq k}}
M_{i+u,j+v}.
\end{aligned}
\label{eq:dilation}
\end{equation}

Similarly, erosion is defined as

\begin{equation}
\begin{aligned}
\bigl[
\operatorname{Erode}(M,k)
\bigr]_{ij}
=
\min_{\substack{-k\leq u\leq k\\-k\leq v\leq k}}
M_{i+u,j+v}.
\end{aligned}
\label{eq:erosion}
\end{equation}
\paragraph{Loss Formulation.}
For an original mask $M$ and its perturbed realization $M'$, we define the changed boundary region by the symmetric difference
\begin{equation}
B(M,M') = \operatorname{Resize}(|M-M'|,h,w),
\label{eq:boundary_band}
\end{equation}
where nearest-neighbor interpolation maps the binary region to latent resolution. The same clean latent $z_0$, Gaussian noise $\epsilon$, and timestep $t$ are used to construct $z_t$ for both conditions. MRDL then compares the two denoising predictions only inside the affected region:
\begin{equation}
\begin{aligned}
\mathcal{L}_{\mathrm{MRDL}}
&=
\mathbb{E}_{M',t,\epsilon}
\left[
\frac{
\left\|
B\odot
\left(
\widehat{\epsilon}_{M}
-
\operatorname{sg}[\widehat{\epsilon}_{M'}]
\right)
\right\|_2^2
}{
\|B\|_1+\eta
}
\right],\\[-2pt]
&\hspace{12mm} B=B(M,M').
\end{aligned}
\label{eq:mrd_loss}
\end{equation}
where $\operatorname{sg}[\cdot]$ denotes stop-gradient and $\eta=10^{-6}$ prevents division by zero. Samples producing an empty changed region are resampled. This normalized consistency loss prevents large masks from dominating the minibatch objective.

\paragraph{Regularization Interpretation.}
MRDL is a local consistency regularizer over plausible annotation-boundary perturbations. It reduces the sensitivity of the conditional score estimate to small changes in the supplied binary contour; it is not an exact marginalization of the data likelihood over physical smoke boundaries.

\paragraph{Total Training Loss.}
The complete training objective combines the standard diffusion loss
with the proposed MRDL:

\begin{equation}
\mathcal{L}_{\mathrm{total}}
=
(1-\omega)\mathcal{L}_{\mathrm{diff}}
+
\omega\mathcal{L}_{\mathrm{MRDL}},
\label{eq:total_loss}
\end{equation}

where $\omega$ controls the strength of the boundary-consistency
regularization. Based on the ablation study in Section~VI-F, we set
$\omega=0.4$ in all reported experiments.







\subsection{Regularization Analysis of MRDL}
\label{sec:theory_mrdl}

Let $f_\theta(M)=\epsilon_\theta(z_t,t,C(M))$ for fixed $z_t$ and $t$. For sufficiently small contour perturbations $\Delta M=M'-M$, a first-order expansion gives
\begin{equation}
f_\theta(M')-f_\theta(M)\approx J_{f_\theta}(M)\Delta M,
\end{equation}
where $J_{f_\theta}(M)$ is the Jacobian of the conditional prediction with respect to the mask condition. Minimizing Eq.~\eqref{eq:mrd_loss} therefore penalizes the locally weighted response $B(M,M')\odot J_{f_\theta}(M)\Delta M$. This encourages local stability around uncertain mask contours while leaving predictions away from the perturbed boundary unconstrained. The analysis motivates MRDL as a Jacobian-style consistency penalty and does not assume that morphological perturbations reproduce the physical evolution of smoke.

\subsection{Multimodal Quality Filtering}

Diffusion stochasticity produces variable-quality outputs. We implement a three-stage filtering pipeline.

\paragraph{Stage 1: Manual Annotation.}
We randomly sample 150 generated images and manually annotate each on three axes (0-10 scale):
\begin{itemize}
    \item \textbf{Color fidelity ($s_{\text{color}}$)}: Alignment with natural smoke tones (grays, browns, blues). Score 10 for perfect color, 0 for completely unnatural colors (e.g., neon green).
    \item \textbf{Visibility ($s_{\text{visibility}}$)}: Smoke prominence against background. Score 10 for clearly visible but not overwhelming, 0 for invisible or fully opaque.
    \item \textbf{Translucency ($s_{\text{translucency}}$)}: Semi-transparency and background blending. Score 10 for realistic transparency gradient, 0 for solid opaque or fully transparent.
\end{itemize}

To ensure consistency, we provide annotators with reference images and detailed rubrics. Inter-annotator agreement measured by Fleiss' $\kappa$ is 0.82, indicating substantial agreement.

\paragraph{VLM Annotation Protocol (Required Definition).}
\label{sec:vlm_protocol}
\textbf{[PENDING VERIFICATION -- every bracketed field below must be
confirmed against the actual annotation records before submission.]}
\begin{itemize}
    \item \textbf{Number and expertise of annotators}: [N] annotators
    with [background/training, e.g., domain familiarity with wildfire
    imagery or a documented calibration session] must be reported.
    \item \textbf{Coverage}: whether all annotators rated all 150 images,
    or a partial/overlapping design was used (which affects how Fleiss'
    $\kappa$ or an ICC should be computed).
    \item \textbf{Disagreement resolution}: the procedure used when
    annotators disagreed (e.g., averaging scores, adjudication by a third
    annotator, or majority vote).
    \item \textbf{Discretization}: Fleiss' $\kappa$ is defined for
    \emph{categorical} ratings. If the underlying 0--10 scores are
    continuous (or effectively so), state explicitly whether they were
    discretized into bins before computing $\kappa=0.82$, and if so, the
    binning scheme used. If the scores were not discretized, an
    intraclass correlation coefficient (ICC, e.g., ICC(2,1) or ICC(3,1)
    depending on whether annotators are treated as a random or fixed
    effect) or another weighted agreement statistic appropriate for
    ordinal/continuous ratings should be reported instead of, or in
    addition to, Fleiss' $\kappa$.
    \item \textbf{Train/validation/test split of the 150 images}: the
    exact number of the 150 human-annotated images used to fit the LoRA
    fine-tuning (Stage 2), to validate hyperparameters, and to hold out
    for reporting prediction error/correlation; this split must be
    disjoint at the image level.
    \item \textbf{Held-out prediction quality}: the fine-tuned VLM's
    prediction error (e.g., MAE/RMSE per score axis) and correlation
    (e.g., Pearson/Spearman) with human ratings on the held-out test
    images, which is currently not reported and is necessary to justify
    using the VLM as a ranking function in Stage 3.
\end{itemize}

\paragraph{Stage 2: VLM Fine-Tuning.}
We fine-tune Qwen2-VL-7B-Instruct \cite{wang2024qwen2vlenhancingvisionlanguagemodels} to predict quality scores. The model receives an image-caption pair and outputs three scores:
\begin{equation}
(s_{\text{color}}, s_{\text{visibility}}, s_{\text{translucency}}) = \text{Qwen2-VL}(x, c)
\end{equation}

Fine-tuning uses LoRA \cite{hu2021loralowrankadaptationlarge} (rank $r=16$) to reduce memory footprint, with learning rate $1\times10^{-4}$ for 10 epochs. The training loss is MSE between predicted and ground-truth scores. Because the annotation set contains 150 images, the VLM is used only as a task-specific ranking component within this pipeline; broader generalization of its absolute scores is not claimed.

\paragraph{Stage 3: Filtering.}
The composite quality metric combines the three scores with empirically determined weights:
\begin{equation}
Q = 0.4 s_{\text{color}} + 0.4 s_{\text{visibility}} + 0.2 s_{\text{translucency}}
\label{eq:quality}
\end{equation}
Translucency receives lower weight because it is inherently more subjective and harder for the VLM to assess reliably.

We generate 96,000 candidate samples, compute $Q$ for each using the fine-tuned VLM, sort the candidates by $Q$, and retain the top 30\% (28,800 samples) for downstream training.

\paragraph{Bounding Box Annotation.}
Since we have the original mask $M$, we can automatically extract bounding boxes without manual labeling. For each mask, we compute:
\begin{equation}
B = \text{BBox}(\text{MaxComponent}(M))
\end{equation}
where $\text{MaxComponent}$ selects the largest connected component (to filter out noise), and $\text{BBox}$ returns $(x_{\min}, y_{\min}, x_{\max}, y_{\max})$. This provides detection labels directly.

\subsection{Distribution Consistency Validation}
\label{sec:distribution_protocol}

A key concern in synthetic data generation is whether generated samples
follow the statistical characteristics of real smoke observations. If the
generated distribution deviates significantly from the real-world
distribution, the resulting detection models may fail to generalize in
operational forest monitoring environments. To address this issue, we
explicitly evaluate the similarity between real and synthetic smoke
datasets using statistical distribution metrics.

\paragraph{Feature Embedding Extraction.}
Let $f_r=\phi(x_r)$ and $f_s=\phi(x_s)$, where $x_r\in D_{\mathrm{real}}$,
$x_s\in D_{\mathrm{synthetic}}$, and $\phi(\cdot)$ denotes the feature
extractor. \textbf{[PENDING VERIFICATION -- inconsistency to resolve]}
Table~\ref{tab:distribution} and Section~\ref{sec:results} state that
these features are extracted with a \emph{CLIP ViT-H/14} image encoder,
whereas this methodology section previously described a ResNet-50
backbone; these two descriptions are inconsistent and must be reconciled
to a single confirmed extractor before submission. Assuming the
CLIP ViT-H/14 extractor (to match Table~\ref{tab:distribution}), the
following must additionally be specified and confirmed against the actual
evaluation code:
\begin{itemize}
    \item \textbf{Exact checkpoint}: [e.g., \texttt{laion/CLIP-ViT-H-14-laion2B-s32B-b79K}
    or the specific OpenCLIP/OpenAI release used] -- must be stated by
    name and source.
    \item \textbf{Feature layer}: [e.g., the pooled image-encoder output
    before the projection head, or the projected embedding] -- state which.
    \item \textbf{Feature normalization}: whether features are
    $\ell_2$-normalized before computing $\mu$, $\Sigma$, FD, and MMD.
    \item \textbf{Number of images}: the exact number of real ($n$) and
    synthetic ($m$) images used to estimate $\mu_r,\Sigma_r$ and
    $\mu_s,\Sigma_s$.
    \item \textbf{Sampling procedure}: how the real and synthetic image
    sets were drawn (e.g., all available real test images vs. a fixed
    random subsample; whether the same synthetic images used for detector
    training are reused here or a held-out synthetic set is used).
\end{itemize}

\paragraph{Distribution Similarity Metrics.}
We compute three complementary metrics.

\emph{Fr\'echet Distance (FD).} We follow the standard Fr\'echet Inception
Distance formulation, applied here to CLIP features rather than Inception
features (hence ``FD'' rather than ``FID''):
\begin{equation}
\mathrm{FD}
=
\|\mu_r-\mu_s\|^2
+
\operatorname{Tr}
\left(
\Sigma_r+\Sigma_s-2(\Sigma_r\Sigma_s)^{1/2}
\right).
\label{eq:fd}
\end{equation}
\textbf{[PENDING VERIFICATION]} The numerical implementation must specify:
(i) the covariance regularization used to stabilize $(\Sigma_r\Sigma_s)^{1/2}$
(e.g., adding $\epsilon I$ with $\epsilon=$~[value] to each covariance
matrix before the matrix square root), and (ii) the specific numerical
routine used to compute the matrix square root (e.g., eigen-decomposition
or the Denman--Beavers iteration, following the standard FID
implementation).

\emph{Maximum Mean Discrepancy (MMD).} Using a kernel $k(\cdot,\cdot)$
(e.g., a Gaussian RBF kernel with bandwidth $\sigma$), the \emph{squared}
MMD between the real and synthetic feature distributions is
\begin{equation}
\begin{aligned}
\mathrm{MMD}^2
={}&
\frac{1}{n^2}\sum_{i,i'=1}^{n}k(f_{r,i},f_{r,i'})
+
\frac{1}{m^2}\sum_{j,j'=1}^{m}k(f_{s,j},f_{s,j'})\\
&-
\frac{2}{nm}\sum_{i=1}^{n}\sum_{j=1}^{m}k(f_{r,i},f_{s,j}).
\end{aligned}
\label{eq:mmd}
\end{equation}
\textbf{[PENDING VERIFICATION]} Table~\ref{tab:distribution} must state
explicitly whether the reported ``MMD'' values are $\mathrm{MMD}^2$ (as in
Eq.~\eqref{eq:mmd}) or $\sqrt{\mathrm{MMD}^2}$, and must specify the kernel
and its bandwidth (or bandwidth-selection procedure, e.g., the median
heuristic).

\emph{Wasserstein Distance.} We report the 2-Wasserstein distance between
the (Gaussian-approximated) real and synthetic feature distributions,
which admits the closed form
\begin{equation}
W_2^2
=
\|\mu_r-\mu_s\|^2
+
\operatorname{Tr}
\left(
\Sigma_r+\Sigma_s-2(\Sigma_r^{1/2}\Sigma_s\Sigma_r^{1/2})^{1/2}
\right),
\label{eq:wasserstein}
\end{equation}
and we report $W_2=\sqrt{W_2^2}$ in Table~\ref{tab:distribution}. Under a
Gaussian assumption Eq.~\eqref{eq:wasserstein} coincides numerically with
Eq.~\eqref{eq:fd} (both reduce to the Fr\'echet distance between Gaussians);
\textbf{[PENDING VERIFICATION]} if the reported Wasserstein values in
Table~\ref{tab:distribution} were instead computed via a different
estimator (e.g., a sliced-Wasserstein approximation, or an empirical
optimal-transport solver on the raw feature samples rather than their
Gaussian fit), the exact estimator and its implementation must be stated
here instead of Eq.~\eqref{eq:wasserstein}.

\paragraph{Uncertainty Quantification.}
\textbf{[PENDING VERIFICATION]} Table~\ref{tab:distribution} currently
reports point estimates only. Confidence intervals or bootstrap standard
errors for FD, MMD, and $W_2$ (e.g., obtained by resampling the real and
synthetic feature sets with replacement) should be reported before
submission so that the differences between methods can be assessed for
statistical, not just numerical, significance.

Lower values of FD, MMD, and $W_2$ indicate stronger alignment between
real and generated distributions.

\paragraph{Interpretation.}
This validation is intended to help confirm that the generated smoke
samples are not only visually plausible but also statistically consistent
with real smoke observations in the chosen feature space, which is
relevant to detector robustness under varying environmental conditions;
the numeric values currently reported in Table~\ref{tab:distribution} are
pending verification against a completed evaluation run using the
protocol above (see TODO\_BEFORE\_SUBMISSION.md).
\subsection{Computational Complexity}
\label{sec:complexity}

The proposed method introduces modest overhead.

\paragraph{Dual Encoder.}

Feature extraction adds $\mathcal{O}(HW)$ complexity.

\paragraph{Joint Cross-Attention.}

\begin{equation}
\text{Standard: } \mathcal{O}(N^2 + NL)
\end{equation}

\begin{equation}
\text{JCA: } \mathcal{O}(N^2 + 2NL)
\end{equation}

resulting in less than 10\% computational increase.

\paragraph{MRDL.}

Morphological perturbations operate in $\mathcal{O}(HW)$ time and are negligible relative to diffusion computation.



\section{Experimental Setup}
\label{sec:experiments}

\subsection{Datasets}

\paragraph{Smoke-Free Backgrounds.}
We curate 8,000 high-resolution images ($1920\times1080$ or higher) from diverse sources:
\begin{itemize}
    \item Forest surveillance cameras (40\%)
    \item Public landscape datasets (30\%)
    \item Aerial drone footage (20\%)
    \item Web-scraped forest scenes (10\%)
\end{itemize}

Environmental diversity is ensured across:
\begin{itemize}
    \item \textbf{Seasons:} Summer (40\%), Autumn (30\%), Winter (20\%), Spring (10\%)
    \item \textbf{Illumination:} Day (60\%), Dusk (25\%), Night (15\%)
    \item \textbf{Weather:} Clear (50\%), Haze (30\%), Fog (20\%)
    \item \textbf{Terrain:} Flat (40\%), Hilly (35\%), Mountainous (25\%)
\end{itemize}

\paragraph{Smoke Mask Extraction.}
We use SAM \cite{kirillov2023segment} with bounding box prompts from existing datasets:
\begin{itemize}
    \item FLAME \cite{SHAMSOSHOARA2021108001}: 40,000 masks
    \item HPWREN \cite{los1998high}: 15,000 masks
    \item SMOKE5K \cite{yan2023transmissionguidedbayesiangenerativemodel}: 5,000 masks
\end{itemize}
Total: 60,000 binary smoke masks with varying shapes, sizes, and densities.

\paragraph{Synthetic Generation.}
For each of the 8,000 backgrounds, three masks are selected from the
60,000-mask pool to form 8{,}000~$\times$~3 = 24,000 (background, mask)
pairs; we then generate 4 variations per pair using \method~with different
random seeds, yielding 96,000 synthetic smoke images. After VLM filtering,
we retain the top 30\% (28,800 samples) for downstream training.

\paragraph{Mask-Selection Procedure (Clarification).}
\textbf{[PENDING VERIFICATION]} The exact procedure used to draw the three
masks assigned to each background must be reported explicitly rather than
left implicit, including:
\begin{itemize}
    \item Whether masks are drawn uniformly at random from the full
    60,000-mask pool, or with constraints (e.g., matching mask aspect ratio
    or spatial extent to the background resolution, or balancing mask
    source dataset -- FLAME/HPWREN/SMOKE5K -- across backgrounds).
    \item Whether sampling is with or without replacement \emph{across}
    backgrounds (a given mask may legitimately be reused on multiple
    backgrounds) and, separately, whether the three masks assigned to the
    \emph{same} background are constrained to be distinct.
    \item Whether the assignment was fixed by a recorded random seed (so
    it is exactly reproducible) or regenerated at each run.
    \item Whether any mask (or the background, or the underlying source
    scene/video the mask was extracted from) that appears in the training
    split also appears in the validation or test split -- see the dataset
    split and leakage-control protocol below, which this mask-assignment
    procedure must respect.
\end{itemize}
This information must be pulled from the actual generation logs/scripts and
inserted here verbatim; it must not be inferred or assumed.

\subsection{Dataset Splits and Leakage Control}
\label{sec:dataset_splits}

\textbf{[PENDING VERIFICATION -- template below; every bracketed field must
be filled from the actual dataset manifests before submission, and no field
may be guessed.]} A complete split specification must report, at minimum:

\begin{itemize}
    \item \textbf{Real detector-training data:} [N\textsubscript{train}]
    training images, [N\textsubscript{val}] validation images, and
    [N\textsubscript{test}] test images, drawn from [dataset name(s), e.g.,
    a named subset of FLAME/HPWREN/SMOKE5K or an independently collected
    real wildfire-smoke detection set]. State whether the official
    published split of each source dataset was used, or whether a custom
    split was created (and if so, the exact split ratios and the random
    seed used to create it).
    \item \textbf{Scene/video-level leakage control:} state explicitly
    whether frames sampled from the same video or camera sequence (e.g.,
    consecutive or temporally adjacent FIgLib/HPWREN frames) are
    constrained to remain within a single split. \emph{Scene-level or
    video-level splitting (rather than frame-level random splitting) must
    be used} to prevent near-duplicate or temporally correlated frames
    from leaking between training and evaluation; the manuscript must
    state which strategy was actually applied.
    \item \textbf{Duplicate/near-duplicate handling:} describe how
    duplicate or temporally adjacent frames were detected (e.g., via
    perceptual hashing, fixed frame-sampling intervals, or manual curation)
    and how many frames were removed as a result.
    \item \textbf{Background/mask/scene overlap between synthetic-data
    splits:} state whether the 8,000 smoke-free backgrounds and 60,000
    masks used to generate the 96,000 synthetic images are themselves
    partitioned into background-level train/val/test splits (so that no
    background used to generate \emph{training} synthetic images also
    appears, or shares a source scene with, the images used for
    validation/testing), or whether backgrounds/masks are shared across
    splits for the synthetic-generation pipeline. If shared, justify why
    this does not constitute leakage for the specific evaluation being
    performed.
    \item \textbf{Composition of the ``fixed base dataset''} referenced in
    Table~\ref{tab:real_vs_synth}: the exact number and source of images it
    contains, and how it was held fixed across the Real+Real and
    Real+Synthetic arms.
    \item \textbf{Value of $N$} in the real-versus-synthetic comparison
    (Table~\ref{tab:real_vs_synth}): the exact number of additional real
    images (Real+Real arm) and additional synthetic images
    (Real+Synthetic arm) added to the fixed base dataset, and confirmation
    that $N$ is identical across both arms for a given model row.
\end{itemize}

\subsection{Implementation Details}

\paragraph{Hardware.}
\begin{itemize}
    \item GPU: NVIDIA A100 80GB (single)
    \item CPU: 32-core Intel Xeon
    \item RAM: 256GB
    \item Storage: 4TB NVMe SSD
\end{itemize}

\paragraph{Training Hyperparameters.}
\begin{table}[H]
\centering
\caption{\textbf{Training hyperparameters.}}
\label{tab:hyperparams}
\begin{tabular}{ll}
\toprule
Parameter & Value \\
\midrule
Optimizer & AdamW ($\beta_1=0.9$, $\beta_2=0.999$) \\
Learning rate & $1\times10^{-4}$ (cosine decay) \\
Warmup steps & 1,000 \\
Weight decay & $1\times10^{-2}$ \\
Batch size & 32 \\
Total iterations & 20,000 \\
Precision & Mixed-precision (fp16) \\
Gradient checkpointing & Enabled \\
EMA decay & 0.9999 \\
\bottomrule
\end{tabular}
\end{table}

\paragraph{Inference Hyperparameters.}
\begin{table}[H]
\centering
\caption{\textbf{Inference hyperparameters.}}
\label{tab:inference_hp}
\begin{tabular}{ll}
\toprule
Parameter & Value \\
\midrule
Sampling algorithm & DDIM \cite{song2022denoisingdiffusionimplicitmodels} \\
Number of steps & 50 \\
Classifier-free guidance scale & 7.5 \\
Latent noise seed (deterministic eval.\ only) & Fixed (42) \\
Output resolution & $512\times512$ \\
\bottomrule
\end{tabular}
\end{table}

\paragraph{Seed Usage Clarification.}
The fixed seed reported above (42) is used only for the deterministic,
reproducible qualitative visualizations and quantitative generative-quality
comparisons in Section~\ref{sec:results} (Tables~\ref{tab:quantitative},
\ref{tab:distribution}, \ref{tab:boundary} and Figures~\ref{fig:ablation_vis},
\ref{fig:qualitative}), so that every compared method is evaluated on
identical noise draws. It is \emph{not} the seed used to generate the
96{,}000-image synthetic training corpus described above, which instead
used four distinct seeds per (background, mask) pair to obtain diverse
variations. \textbf{[PENDING VERIFICATION]} the exact seed sequence used for
corpus generation (e.g., a recorded list such as
$\{0,1,2,3\}$ or logged per-sample seeds from the generation run) must be
retrieved from the generation logs and reported here verbatim before
submission; it must not be re-derived or guessed.

\paragraph{Architecture Details.}
\begin{table}[H]
\centering
\caption{U-Net architecture details.}
\label{tab:unet}
\resizebox{0.98\linewidth}{!}{%
\begin{tabular}{lcccc}
\toprule
Stage & Resolution & Channels & Attention & Conditioning \\
\midrule
Input      & $64\times64$ & 320  & --            & --    \\
Down 1     & $32\times32$ & 320  & --            & $F_m$ \\
Down 2     & $16\times16$ & 640  & --            & $F_m$ \\
Down 3     & $8\times8$   & 1280 & $\checkmark$  & JCA   \\
Bottleneck & $8\times8$   & 1280 & $\checkmark$  & $F_m$ \\
Up 1       & $8\times8$   & 1280 & $\checkmark$  & JCA   \\
Up 2       & $16\times16$ & 640  & --            & $F_M$ \\
Up 3       & $32\times32$ & 320  & --            & $F_M$ \\
Up 4       & $64\times64$ & 320  & --            & $F_M$ \\
Output     & $64\times64$ & 320  & --            & --    \\
\bottomrule
\end{tabular}%
}
\end{table}

\subsection{Detector Training Configuration}
\label{sec:detector_config}

Table~\ref{tab:hyperparams} and Table~\ref{tab:inference_hp} describe the
training and sampling configuration of the \emph{diffusion generator},
not of the downstream YOLO detectors evaluated in Section~\ref{sec:results}.
Table~\ref{tab:detector_config} below is a dedicated Detector Training
Configuration table.
\textbf{[PENDING VERIFICATION -- every bracketed field must be filled from
the actual detector training logs/configs before submission; no field may
be guessed.]}

\begin{table}[t]
\centering
\caption{\textbf{Detector training configuration (template).} One column
must be completed per YOLO variant (v6--v13) evaluated in
Table~\ref{tab:detection}; only the shared/common fields are shown as a
single column here for brevity, but per-variant values (e.g., differing
batch size or learning rate across scales) must be reported if they
differ.}
\label{tab:detector_config}
\begin{tabular}{ll}
\toprule
Parameter & Value \\
\midrule
Model variant / scale (per version) & [e.g., YOLOv13-N/S/L/X -- state which] \\
Repository commit / release tag & [exact commit hash or release tag] \\
Input resolution & [value, e.g., 640$\times$640] \\
Batch size & [value] \\
Optimizer & [e.g., SGD / AdamW] \\
Initial learning rate \& schedule & [value and schedule] \\
Number of epochs & [value] \\
Pretrained weights source & [e.g., COCO-pretrained checkpoint, source/URL] \\
Data augmentation settings & [e.g., Mosaic, Mixup, HSV jitter -- exact config] \\
Confidence threshold & [value] \\
NMS IoU threshold & [value] \\
Early-stopping policy & [patience / criterion, or ``none''] \\
Random seed(s) & [list of seeds used per run] \\
Hardware & [GPU model, count] \\
Training wall-clock time & [value per run] \\
\bottomrule
\end{tabular}
\end{table}

\paragraph{YOLOv12 and YOLOv13 Configurations.}
YOLOv12 was implemented using [release/commit] from the official
repository \cite{tian2025yolov12}, using the [exact model variant, e.g.,
YOLOv12-S]. YOLOv13 was implemented using [release/commit] from the
official repository \cite{lei2025yolov13}, using the [exact model
variant, e.g., YOLOv13-L]. Both models were trained at an input
resolution of [value] for [value] epochs with batch size [value]. We used
[optimizer], an initial learning rate of [value], and pretrained weights
from [source]. Confidence and NMS thresholds were set to [values]. All
comparisons used identical train/validation/test splits (Section~\ref{sec:dataset_splits})
and random seeds [list]. \textbf{[PENDING VERIFICATION]} every bracketed
field above must be replaced with the actual configuration used, sourced
from the training logs/config files for the reported YOLOv12/YOLOv13 runs;
it must not be guessed or inferred from the official repository defaults
without confirming that those defaults were in fact used.

\subsection{Evaluation Metrics}

\paragraph{Generative Quality Metrics.}
\begin{itemize}
    \item \textbf{PSNR (dB)}: Peak signal-to-noise ratio. Higher is better.
    \item \textbf{SSIM} \cite{1284395}: Structural similarity index. Range [0,1], higher is better.
    \item \textbf{LPIPS} \cite{zhang2018unreasonableeffectivenessdeepfeatures}: Learned perceptual similarity. Lower is better; we report 1-LPIPS for intuitive comparison.
    \item \textbf{CLIP-Sim} \cite{radford2021learningtransferablevisualmodels}: Cosine similarity between image and text prompt embeddings. Higher is better.
\end{itemize}

\paragraph{PSNR/SSIM/LPIPS Comparison Protocol (Required Definition).}
\label{sec:psnr_protocol}
PSNR, SSIM, and LPIPS are only scientifically interpretable when the exact
pair of images being compared, and the pixels included in that comparison,
are unambiguous. \textbf{[PENDING VERIFICATION -- the bracketed choices
below must be confirmed against the actual evaluation code before
submission; they must not be asserted without that confirmation.]} The
protocol must specify:
\begin{itemize}
    \item \textbf{Reference target}: whether the generated image is
    compared against (a) a paired real target image depicting the same
    scene with real smoke, (b) the original smoke-free background image
    restricted to the region \emph{outside} the mask (i.e., a
    preserved-background-fidelity check), or (c) some other reference.
    \method~does not have access to a real paired ``ground-truth smoke''
    target for a given synthetic mask, so option (b) -- background
    preservation outside $M$ -- is the only protocol that yields a
    well-defined pixel-wise reference; this must be stated explicitly as
    [confirmed protocol].
    \item \textbf{Region of comparison}: whether the metric is computed
    over the whole image, only the preserved (unmasked) region, or only
    the mask region $M$ (in which case PSNR/SSIM against a background
    reference are not meaningful, since the mask region is intentionally
    changed by design).
    \item \textbf{Spatial alignment}: how generated and reference images
    are spatially aligned (e.g., identical background pixel grid by
    construction, since inpainting preserves unmasked pixel coordinates
    exactly) and whether any resizing/cropping is applied before metric
    computation.
    \item \textbf{Pixel value range and color space} used for PSNR/SSIM
    (e.g., $[0,1]$ or $[0,255]$, RGB vs. luminance-only).
\end{itemize}

\paragraph{Baseline Reproducibility.}
\label{sec:baseline_repro}
Table~\ref{tab:baseline_repro} specifies the reproducibility details
required for every baseline compared against \method~in
Table~\ref{tab:quantitative} and Table~\ref{tab:sources}.
\textbf{[PENDING VERIFICATION -- template below; every bracketed field
must be filled from the actual baseline evaluation runs before
submission.]}

\begin{table*}[t]
\centering
\caption{\textbf{Baseline reproducibility specification (template).} All
baselines must be evaluated under an equal sample-count budget for a fair
comparison; the ``Equal-budget?'' column must confirm this explicitly.}
\label{tab:baseline_repro}
\small
\begin{tabular}{lccccccccc}
\toprule
Baseline & Checkpoint/version & Prompt template & Mask preprocessing & Sampling method/steps & Guidance scale & Output res. & Fine-tuned? & \# samples & Filtering applied \\
\midrule
Stable Diffusion \cite{rombach2022highresolutionimagesynthesislatent} & [ckpt] & [template] & [preproc.] & [method/steps] & [value] & [res.] & [yes/no] & [N] & [filter/none] \\
PowerPaint \cite{zhuang2024taskworthwordlearning} & [ckpt] & [template] & [preproc.] & [method/steps] & [value] & [res.] & [yes/no] & [N] & [filter/none] \\
SD-Inpainting & [ckpt] & [template] & [preproc.] & [method/steps] & [value] & [res.] & [yes/no] & [N] & [filter/none] \\
BLD \cite{Avrahami_2023} & [ckpt] & [template] & [preproc.] & [method/steps] & [value] & [res.] & [yes/no] & [N] & [filter/none] \\
FlameDiffuser \cite{wang2024flamediffuserwildfireimage} & [ckpt] & [template] & [preproc.] & [method/steps] & [value] & [res.] & [yes/no] & [N] & [filter/none] \\
GAN-based synthesis \cite{rs12223715} & [ckpt] & [n/a] & [preproc.] & [n/a] & [n/a] & [res.] & [yes/no] & [N] & [filter/none] \\
\textbf{\method{} (Ours)} & [ckpt] & [template] & [preproc.] & [method/steps] & [value] & [res.] & yes & [N] & VLM top-30\% \\
\bottomrule
\end{tabular}
\end{table*}

Without the information in Table~\ref{tab:baseline_repro} populated from
the actual baseline runs, the comparisons in
Table~\ref{tab:quantitative}/Table~\ref{tab:sources} cannot currently be
confirmed as controlled (equal sample budget, comparable guidance scale,
comparable fine-tuning status), and must be treated as
provisional/illustrative until this table is completed.

\paragraph{Detection Performance Metrics.}
\begin{itemize}
    \item \textbf{mAP@50}: Mean average precision at IoU threshold 0.5.
    \item \textbf{mAP@50--95}: Average over IoU thresholds 0.5:0.05:0.95.
    \item \textbf{Precision}: $\text{TP}/(\text{TP}+\text{FP})$.
    \item \textbf{Recall}: $\text{TP}/(\text{TP}+\text{FN})$.
\end{itemize}

\paragraph{Boundary Quality Metrics.}
\label{sec:boundary_metric_protocol}
The boundary band is defined as
\begin{equation}
B_{\delta}(M)
=
\operatorname{Dilate}(M,\delta)
\setminus
\operatorname{Erode}(M,\delta),
\qquad \delta=5.
\label{eq:boundary_band_metric}
\end{equation}

The boundary-softness score is computed as
\begin{equation}
\mathcal{S}
=
\frac{1}{|B_{\delta}|}
\sum_{p\in B_{\delta}}
\left\lVert\nabla I(p)\right\rVert_2 .
\label{eq:boundary_softness}
\end{equation}

The gradient-distribution divergence is
\begin{equation}
\mathcal{G}
=
D_{\mathrm{KL}}
\left(
p_{\mathrm{real}}(g)
\,\middle\|\,
p_{\mathrm{synth}}(g)
\right),
\label{eq:gradient_divergence}
\end{equation}
where
\begin{equation}
g(p)=\left\lVert\nabla I(p)\right\rVert_2,
\qquad p\in B_{\delta}.
\label{eq:gradient_magnitude}
\end{equation}

\paragraph{Boundary-Metric Implementation Details (Required Definition).}
\textbf{[PENDING VERIFICATION -- the bracketed choices below must be
confirmed against the actual evaluation code before submission.]} A lower
gradient magnitude $\mathcal{S}$ is not unconditionally better, since
excessive blurring can also depress $\mathcal{S}$ without indicating a
more realistic boundary; $\mathcal{S}$ must therefore always be reported
as a \emph{distance from the real-smoke reference value} (as in
Table~\ref{tab:boundary}), not interpreted in isolation. The full protocol
must additionally specify:
\begin{itemize}
    \item \textbf{Color space / grayscale conversion}: whether $\nabla I$
    is computed on [grayscale-converted luminance / each RGB channel
    averaged / a specific channel], and the exact conversion formula used.
    \item \textbf{Gradient operator}: whether $\nabla I$ is computed with
    a [Sobel / Scharr / central finite-difference] operator, and the
    kernel size used.
    \item \textbf{Pixel-value normalization}: the intensity range (e.g.,
    $[0,1]$ or $[0,255]$) used before differentiation.
    \item \textbf{Boundary-band handling near image borders}: how
    $B_{\delta}(M)$ and $\nabla I$ are evaluated when the dilated/eroded
    mask or the gradient stencil extends past the image boundary (e.g.,
    zero-padding, reflection, or exclusion of border pixels).
    \item \textbf{Histogram binning and smoothing for $\mathcal{G}$}: the
    number of histogram bins used to estimate $p_{\mathrm{real}}(g)$ and
    $p_{\mathrm{synth}}(g)$ in Eq.~\eqref{eq:gradient_divergence}, any
    smoothing/regularization applied to avoid zero-probability bins, and
    whether $g$ is binned on a fixed or data-dependent range.
    \item \textbf{Sample size and averaging}: the number of images
    averaged to obtain $\mathcal{S}$ and $\mathcal{G}$ per method, and
    whether $\mathcal{S}$ is averaged per-image-then-across-images or
    pooled across all boundary pixels from all images.
    \item \textbf{Uncertainty estimates}: standard deviation, standard
    error, or bootstrap confidence intervals for $\mathcal{S}$ across the
    evaluated images, which are currently not reported in
    Table~\ref{tab:boundary}.
\end{itemize}

\section{Results}
\label{sec:results}

\textbf{Editorial notice on result verification status.} \textbf{[PENDING
VERIFICATION -- applies to every table and figure in this section]} Per
the reviewer assessment accompanying this revision, every numeric value
reported in this section (Tables~\ref{tab:quantitative}--\ref{tab:filter_equal_size},
Figures~\ref{fig:mrdl_sensitivity}, \ref{fig:ablation_vis},
\ref{fig:qualitative}, the FD/MMD/Wasserstein values in
Table~\ref{tab:distribution}, the boundary-softness/KL values in
Table~\ref{tab:boundary}, the five-run means/standard deviations in
Table~\ref{tab:domainshift}, and the reported $p<0.01$ significance claim)
must be traced to a completed experiment with retained experiment logs,
evaluation scripts, and saved model predictions before this manuscript is
submitted. This section is presented as a \emph{proposed evaluation
protocol with illustrative placeholder values} rather than as verified
experimental findings until that verification is complete; see
TODO\_BEFORE\_SUBMISSION.md for the corresponding checklist. Values must
not be assumed, estimated, or presented as final results in any
submitted version of this manuscript.

\subsection{Quantitative Generative Quality}

Table~\ref{tab:quantitative} compares the generative performance of the evaluated methods. \method~achieves the strongest reported values for PSNR, SSIM, 1-LPIPS, and CLIP similarity:

\begin{itemize}
    \item \textbf{PSNR:} 28.10 dB, an 8.1\% improvement over SD-Inpainting (26.00 dB)
    \item \textbf{SSIM:} 0.88, reflecting superior structural preservation
    \item \textbf{1-LPIPS:} 0.92, indicating the lowest perceptual distance among the evaluated methods
    \item \textbf{CLIP-Sim:} 0.258, demonstrating strong text-image alignment
\end{itemize}

\begin{table}[t]
\centering
\caption{Quantitative comparison of generative performance. Best results are shown in bold.}
\label{tab:quantitative}
\resizebox{\columnwidth}{!}{%
\begin{tabular}{lcccc}
\toprule
Method & PSNR $\uparrow$ & SSIM $\uparrow$ & 1-LPIPS $\uparrow$ & CLIP-Sim $\uparrow$ \\
\midrule
Stable Diffusion \cite{rombach2022highresolutionimagesynthesislatent} & 25.10 & 0.78 & 0.87 & 0.251 \\
PowerPaint \cite{zhuang2024taskworthwordlearning} & 25.60 & 0.80 & 0.89 & 0.253 \\
SD-Inpainting & 26.00 & 0.81 & 0.90 & 0.255 \\
BLD \cite{Avrahami_2023} & 26.30 & 0.83 & 0.91 & 0.254 \\
FlameDiffuser \cite{wang2024flamediffuserwildfireimage} & 27.10 & 0.85 & 0.91 & 0.256 \\
\textbf{\method{} (Ours)} & \textbf{28.10} & \textbf{0.88} & \textbf{0.92} & \textbf{0.258} \\
\bottomrule
\end{tabular}%
}
\end{table}

\subsection{Distribution Alignment}

Table~\ref{tab:distribution} evaluates distributional similarity between real and synthetic smoke in feature space. Using CLIP ViT-H/14 features, \method~achieves the lowest FD (28.15), MMD (0.062), and Wasserstein distance (5.14), confirming that our synthetic data lies closer to the real distribution than prior methods by substantial margins (28\% lower FD than FlameDiffuser).

\begin{table}[t]
\centering
\caption{\textbf{Distribution alignment between real and synthetic smoke.} Lower is better. Features from CLIP ViT-H/14.}
\label{tab:distribution}
\begin{tabular}{lccc}
\toprule
Method & FD$\downarrow$ & MMD$\downarrow$ & Wasserstein$\downarrow$ \\
\midrule
GAN-based synthesis \cite{rs12223715} & 45.32 & 0.124 & 8.67 \\
Vanilla SD-inpainting & 38.47 & 0.098 & 7.23 \\
FlameDiffuser \cite{wang2024flamediffuserwildfireimage} & 32.15 & 0.081 & 6.18 \\
\textbf{\method{} (Ours)} & \textbf{28.15} & \textbf{0.062} & \textbf{5.14} \\
\bottomrule
\end{tabular}
\end{table}

\subsection{Boundary Softness Analysis}

Table~\ref{tab:boundary} quantifies boundary quality. \method~produces gradient distributions closest to the real-smoke reference among the evaluated methods (KL divergence 0.041 vs. SD-Inpainting's 0.187); this is a lower divergence than the baselines, not evidence of distributional identity, and is reported alongside the boundary-softness measure below because a lower gradient magnitude alone can also result from excessive blurring rather than realistic diffuseness (see the boundary-metric protocol in Section~\ref{sec:boundary_metric_protocol}). The boundary softness $\mathcal{S}$ is also closest to the real-smoke reference among the evaluated methods (0.124 vs. 0.119 for real).

\begin{table}[t]
\centering
\caption{\textbf{Boundary quality metrics.} Real smoke values shown for reference.}
\label{tab:boundary}
\begin{tabular}{lcc}
\toprule
Method & Boundary Softness $\mathcal{S}$ & KL Divergence $\downarrow$ \\
\midrule
Real smoke & 0.119 & - \\
SD-Inpainting & 0.203 & 0.187 \\
FlameDiffuser & 0.167 & 0.112 \\
\textbf{\method} & \textbf{0.124} & \textbf{0.041} \\
\bottomrule
\end{tabular}
\end{table}

\subsection{Downstream Detection Performance}

\paragraph{Main Results.}
Table~\ref{tab:detection} reports detection metrics across YOLO variants (v6-v13). Training with \method~synthetic data yields consistent and substantial improvements:

\begin{itemize}
    \item \textbf{YOLOv13:} mAP@50 improves from 0.683 $\rightarrow$ 0.829 (+21.4\% relative)
    \item \textbf{YOLOv10:} Recall improves from 0.593 $\rightarrow$ 0.791 (+33.4\% relative)
    \item \textbf{Average across all models:} +0.156 mAP@50, +0.126 Recall
\end{itemize}

\begin{table*}[t]
\centering
\caption{\textbf{Performance comparison of smoke detection models trained on real vs. mixed datasets.} Mixed dataset combines real data with \method~synthetic samples.}
\label{tab:detection}
\small
\begin{tabular}{lcccccccc}
\toprule
& \multicolumn{4}{c}{Real Dataset Only} & \multicolumn{4}{c}{Mixed Dataset (Real + \method)} \\
\cmidrule(lr){2-5} \cmidrule(lr){6-9}
Model & mAP@50 & mAP@50--95 & Precision & Recall & mAP@50 & mAP@50--95 & Precision & Recall \\
\midrule
YOLOv6 & 0.582 & 0.285 & 0.668 & 0.552 & 0.752 & 0.362 & 0.724 & 0.676 \\
YOLOv7 & 0.598 & 0.291 & 0.677 & 0.568 & 0.761 & 0.371 & 0.735 & 0.688 \\
YOLOv8 & 0.621 & 0.309 & 0.689 & 0.612 & 0.789 & 0.399 & 0.765 & 0.719 \\
YOLOv9 & 0.645 & 0.323 & 0.705 & 0.625 & 0.791 & 0.409 & 0.778 & 0.728 \\
YOLOv10 & 0.603 & 0.294 & 0.672 & 0.593 & 0.776 & 0.384 & 0.749 & 0.791 \\
YOLOv11 & 0.648 & 0.334 & 0.756 & 0.604 & 0.794 & 0.414 & 0.802 & 0.729 \\
YOLOv12 & 0.672 & 0.348 & 0.778 & 0.638 & 0.813 & 0.448 & 0.831 & 0.751 \\
YOLOv13 & 0.683 & 0.359 & 0.791 & 0.649 & \textbf{0.829} & \textbf{0.531} & \textbf{0.847} & \textbf{0.763} \\
\bottomrule
\end{tabular}
\end{table*}

\paragraph{Comparison with Real Data Augmentation.}
Table~\ref{tab:real_vs_synth} compares synthetic augmentation against adding more real data (both settings add $N$ samples to a fixed base dataset). Remarkably, \method~synthetic data \emph{outperforms} additional real samples across all metrics, validating the utility of targeted synthetic generation. For YOLOv13, synthetic data achieves mAP@50 of 0.829 vs. 0.701 for additional real data—an 18.3\% advantage.

\begin{table}[H]
\centering
\caption{Real versus synthetic data augmentation comparison.}
\label{tab:real_vs_synth}
\resizebox{0.98\linewidth}{!}{%
\begin{tabular}{lcccccccc}
\toprule
& \multicolumn{4}{c}{Real + Real}
& \multicolumn{4}{c}{Real + Synthetic} \\
\cmidrule(lr){2-5}
\cmidrule(lr){6-9}

Model
& \shortstack{mAP\\@50}
& \shortstack{mAP\\@50--95}
& Prec.
& Rec.
& \shortstack{mAP\\@50}
& \shortstack{mAP\\@50--95}
& Prec.
& Rec. \\
\midrule

YOLOv6
& 0.637 & 0.330 & 0.641 & 0.642
& 0.752 & 0.362 & 0.724 & 0.676 \\

YOLOv8
& 0.675 & 0.350 & 0.710 & 0.665
& 0.789 & 0.399 & 0.765 & 0.719 \\

YOLOv10
& 0.670 & 0.341 & 0.737 & 0.636
& 0.776 & 0.384 & 0.749 & \textbf{0.791} \\

YOLOv13
& \textbf{0.701}
& \textbf{0.362}
& \textbf{0.779}
& \textbf{0.682}
& \textbf{0.829}
& \textbf{0.531}
& \textbf{0.847}
& 0.763 \\
\bottomrule
\end{tabular}%
}
\end{table}

\paragraph{Comparison with Other Synthetic Sources.}
Table~\ref{tab:sources} compares different synthetic data sources under identical training budgets. \method~substantially outperforms GAN-based and vanilla diffusion baselines, achieving mAP@50 of 0.829 vs. 0.712 (GAN) and 0.748 (vanilla diffusion).

\begin{table}[t]
\centering
\caption{\textbf{Downstream detection performance using different synthetic data sources.} YOLOv13 detector.}
\label{tab:sources}
\begin{tabular}{lcccc}
\toprule
Training Data & mAP@50 & mAP@50--95 & Precision & Recall \\
\midrule
Real only & 0.683 & 0.359 & 0.791 & 0.649 \\
Real + GAN synth & 0.712 & 0.382 & 0.805 & 0.671 \\
Real + Vanilla diffusion synth & 0.748 & 0.412 & 0.821 & 0.698 \\
\textbf{Real + \method{} synthetic} & \textbf{0.829} & \textbf{0.531} & \textbf{0.847} & \textbf{0.763} \\
\bottomrule
\end{tabular}
\end{table}

\subsection{Domain Generalization Under Environmental Shifts}

Table~\ref{tab:domainshift} evaluates cross-domain robustness. \method~synthetic data substantially improves generalization across all three shift scenarios:

\begin{itemize}
    \item \textbf{Season shift (Summer $\rightarrow$ Winter):} mAP@50 improves from 0.412 $\rightarrow$ 0.587 (+42.5\% relative)
    \item \textbf{Illumination shift (Day $\rightarrow$ Dusk):} mAP@50 improves from 0.385 $\rightarrow$ 0.542 (+40.8\%)
    \item \textbf{Weather shift (Clear $\rightarrow$ Haze/Fog):} mAP@50 improves from 0.338 $\rightarrow$ 0.496 (+46.7\%)
\end{itemize}

The largest improvement occurs under weather shift, where background ambiguity is strongest—confirming that boundary-aware diffusion synthesis improves robustness in low-contrast conditions. \textbf{[PENDING VERIFICATION -- see statistical testing protocol below]} we currently report ``all improvements are statistically significant ($p<0.01$)''; this claim requires the complete protocol described next before it can be confirmed.

\paragraph{Statistical Testing Protocol (Required Definition).}
\label{sec:stats_protocol}
With only 5 runs per condition, a paired $t$-test's normality assumption is
difficult to justify from the data alone, and the following must be
reported explicitly rather than a single aggregate $p$-value:
\begin{itemize}
    \item \textbf{Paired observation definition}: each pair consists of
    the (Real Only, Real+\method) mAP@50 (or Recall) obtained from the
    \emph{same} random seed and the \emph{same} train/test split for a
    given domain-shift row, so that the pairing controls for
    run-to-run variation due to seed and split.
    \item \textbf{Five random seeds}: [PENDING] the explicit list of the
    five seeds used for the paired runs (e.g., $\{0,1,2,3,4\}$ or the
    logged seed values) must be reported, not merely ``5 runs.''
    \item \textbf{Test choice and assumptions}: given $n=5$ paired
    differences per row, a paired permutation test or a Wilcoxon
    signed-rank test is likely more defensible than a paired $t$-test,
    which assumes approximately normally distributed differences -- a
    strong assumption with only 5 samples. If a $t$-test is retained, a
    normality check (e.g., Shapiro--Wilk on the 5 paired differences, or
    an explicit justification for why normality is assumed) must be
    reported.
    \item \textbf{Exact $p$-values}: report the specific $p$-value for
    each of the three domain-shift rows (and for mAP@50 and Recall
    separately), rather than a single blanket ``$p<0.01$'' for all
    comparisons.
    \item \textbf{Effect sizes and confidence intervals}: report an
    effect size (e.g., Cohen's $d_z$ for paired differences, or the
    Hodges--Lehmann estimator if using Wilcoxon) together with a
    confidence interval for the mean/median improvement in each row.
    \item \textbf{Multiple-comparison correction}: since 3 domain-shift
    rows $\times$ 2 metrics = 6 hypothesis tests are performed, a
    correction (e.g., Holm--Bonferroni or Benjamini--Hochberg) must be
    applied and the corrected significance threshold reported.
\end{itemize}
Table~\ref{tab:domainshift} below reports mean $\pm$ standard deviation
over 5 runs; \textbf{[PENDING VERIFICATION]} the underlying per-seed
values and the statistical test outputs described above must be reported
(e.g., in a supplementary table) before submission.

\begin{table}[t]
\centering
\caption{\textbf{Generalization under monitoring domain shifts.} Mean $\pm$ std over 5 runs. YOLOv13 detector.}
\label{tab:domainshift}
\begin{tabular}{lcccc}
\toprule
& \multicolumn{2}{c}{Real Only} & \multicolumn{2}{c}{Real + \method} \\
\cmidrule(lr){2-3} \cmidrule(lr){4-5}
Train $\rightarrow$ Test & mAP@50 & Recall & mAP@50 & Recall \\
\midrule
Summer $\rightarrow$ Winter & 0.412$\pm$0.023 & 0.385$\pm$0.021 & 0.587$\pm$0.018 & 0.542$\pm$0.019 \\
Day $\rightarrow$ Dusk & 0.385$\pm$0.031 & 0.352$\pm$0.027 & 0.542$\pm$0.022 & 0.498$\pm$0.020 \\
Clear $\rightarrow$ Haze/Fog & 0.338$\pm$0.028 & 0.306$\pm$0.024 & 0.496$\pm$0.025 & 0.458$\pm$0.021 \\
\midrule
\textbf{Average} & 0.378 & 0.348 & \textbf{0.542} & \textbf{0.499} \\
\bottomrule
\end{tabular}
\end{table}

\FloatBarrier

\subsection{Ablation Studies}

\paragraph{Component Ablation.}
Table~\ref{tab:ablation} ablates key architectural components. The full model outperforms all variants. Percentages below are \emph{relative} decreases computed with the full-model score as the denominator (e.g., $(0.829-0.758)/0.829\times100=8.6\%$); the corresponding absolute change in percentage points is also given for clarity:
\begin{itemize}
    \item Removing JCA reduces mAP@50 from 0.829 $\rightarrow$ 0.758 ($-$0.071 percentage points, $-$8.6\% relative)
    \item Removing MRDL reduces mAP@50 from 0.829 $\rightarrow$ 0.771 ($-$0.058 percentage points, $-$7.0\% relative)
    \item Removing ResNet50 (random initialization) reduces mAP@50 from 0.829 $\rightarrow$ 0.735 ($-$0.094 percentage points, $-$11.3\% relative)
    \item Removing filtering reduces mAP@50 from 0.829 $\rightarrow$ 0.748 ($-$0.081 percentage points, $-$9.8\% relative to the full model; equivalently a $+$10.8\% relative improvement of filtering over random selection when the no-filtering score is used as the denominator, as in Eq.~\eqref{eq:filtering_relative_improvement})
\end{itemize}

\begin{table*}[t]
\centering
\caption{\textbf{Ablation of architectural components.} YOLOv13 detector.
% Rendered as a double-column-spanning table* (rather than a
% single-column table) because the "Configuration" labels (e.g., "No
% ResNet50 (random init)") previously wrapped inside a narrow single
% column, causing this table to visually overlap with the adjacent
% Table~\ref{tab:sources} placed in the neighboring column -- the exact
% rendering defect flagged in the reviewer assessment.
}
\label{tab:ablation}
\begin{tabular}{lcccc}
\toprule
Configuration & mAP@50 & mAP@50--95 & Precision & Recall \\
\midrule
No JCA (CA only) & 0.758 & 0.412 & 0.781 & 0.702 \\
No MRDL ($\omega=0$) & 0.771 & 0.428 & 0.795 & 0.715 \\
No ResNet50 (random init) & 0.735 & 0.389 & 0.748 & 0.681 \\
No filtering (random selection) & 0.748 & 0.412 & 0.761 & 0.698 \\
\textbf{Full \method} & \textbf{0.829} & \textbf{0.531} & \textbf{0.847} & \textbf{0.763} \\
\bottomrule
\end{tabular}
\end{table*}

\FloatBarrier

\paragraph{MRDL Weight Sensitivity.}
Figure~\ref{fig:mrdl_sensitivity} shows mAP@50 as a function of MRDL weight $\omega$. Performance peaks at $\omega=0.4$, with degradation on either side:
\begin{itemize}
    \item $\omega=0$ (no MRDL): mAP@50 = 0.771 (hard edges)
    \item $\omega=0.4$ (optimal): mAP@50 = 0.829
    \item $\omega=1.0$ (MRDL only): mAP@50 = 0.682 (loss of spatial structure)
\end{itemize}


\begin{figure}[t]
\centering
\includegraphics[width=0.95\columnwidth]{fig31.png}
\caption{\textbf{MRDL weight sensitivity analysis.} Performance peaks at $\omega=0.4$. Too low ($\omega=0$): hard boundary artifacts. Too high ($\omega=1.0$): loss of spatial coherence.}
\label{fig:mrdl_sensitivity}
\end{figure}

\FloatBarrier




\paragraph{Filtering Threshold Sensitivity (Confounded by Dataset Size).}
\label{sec:filtering_confound}
Table~\ref{tab:filter_threshold} reports an exploratory retention-threshold analysis. \textbf{[PENDING VERIFICATION -- see TODO\_BEFORE\_SUBMISSION.md]} The 30\% setting obtains the strongest result in this sweep. However, because the retained dataset size changes simultaneously with the quality threshold (from 9,600 to 96,000 images), this experiment cannot by itself establish that a 30\% retention threshold is optimal, or that VLM-based ranking is responsible for the observed differences rather than dataset size alone. The percentage improvement attributed to filtering in Sections~\ref{sec:results} and~\ref{sec:conclusion} is computed relative to the unfiltered (100\%-retention) score at the \emph{same} downstream evaluation:
\begin{equation}
\Delta_{\mathrm{rel}}
=
\frac{Q_{\mathrm{filtered}}-Q_{\mathrm{unfiltered}}}{Q_{\mathrm{unfiltered}}}\times100,
\label{eq:filtering_relative_improvement}
\end{equation}
where $Q$ denotes mAP@50. This does not, however, resolve the size confound described above.

\emph{Proposed equal-size control (pending execution).} To isolate the contribution of VLM-based ranking from dataset size, we propose comparing four \emph{equal-size} (28,800-sample) synthetic subsets drawn from the same 96,000-image candidate pool:
\begin{itemize}
    \item \textbf{Top-ranked 28,800}: the current top-30\% selection (highest $Q$ from Eq.~\eqref{eq:quality}).
    \item \textbf{Random 28,800}: uniformly sampled without regard to $Q$.
    \item \textbf{Bottom-ranked 28,800}: the lowest-$Q$ samples.
    \item \textbf{Stratified 28,800}: sampled to match the overall $Q$-score distribution of the full 96,000-image pool.
\end{itemize}
Table~\ref{tab:filter_equal_size} reports the template for this control; every cell is a placeholder pending the corresponding detector-training run and must not be populated with assumed values.

\begin{table}[t]
\centering
\caption{\textbf{Proposed equal-size ($N=28{,}800$) filtering control (template -- not yet executed).} All cells require a completed YOLOv13 training run before this table can be reported as a result.}
\label{tab:filter_equal_size}
\begin{tabular}{lccc}
\toprule
Selection Strategy & Dataset Size & mAP@50 & Recall \\
\midrule
Top-ranked (VLM) & 28,800 & [PENDING] & [PENDING] \\
Random & 28,800 & [PENDING] & [PENDING] \\
Bottom-ranked (VLM) & 28,800 & [PENDING] & [PENDING] \\
Stratified & 28,800 & [PENDING] & [PENDING] \\
\bottomrule
\end{tabular}
\end{table}

\begin{table}[t]
\centering
\caption{\textbf{Effect of filtering threshold on detection performance.} YOLOv13 detector. \textbf{[PENDING VERIFICATION]} -- values below must be traced to saved detector checkpoints, evaluation logs, and predictions before submission; see TODO\_BEFORE\_SUBMISSION.md.}
\label{tab:filter_threshold}
\begin{tabular}{lccc}
\toprule
Retention Threshold & Dataset Size & mAP@50 & Recall \\
\midrule
10\% & 9,600 & 0.812 & 0.741 \\
20\% & 19,200 & 0.821 & 0.752 \\
30\% & 28,800 & \textbf{0.829} & \textbf{0.763} \\
40\% & 38,400 & 0.818 & 0.751 \\
50\% & 48,000 & 0.805 & 0.738 \\
No filtering & 96,000 & 0.748 & 0.698 \\
\bottomrule
\end{tabular}
\end{table}

\FloatBarrier

\paragraph{Qualitative Ablation Visualizations.}
Figure~\ref{fig:ablation_vis} shows representative outputs from ablated configurations:
\begin{itemize}
    \item No JCA: Smoke exhibits unrealistic texture repetition and poor background integration.
    \item No MRDL: Sharp, binary boundaries with no semi-transparency.
    \item No ResNet50: Severe artifacts, often failing to produce coherent smoke structures.
    \item No filtering: Occasional catastrophic failures (e.g., smoke in wrong colors).
    \item Full model: Realistic, semi-transparent smoke with natural boundaries.
\end{itemize}

\begin{figure}[t]
\centering
\IfFileExists{qualitative_ablation.pdf}{%
\includegraphics[width=0.98\linewidth]{qualitative_ablation.pdf}%
}{\PackageError{FireSmokeGenNet}{Missing qualitative_ablation.pdf}{Upload the experimentally generated qualitative ablation panel before submission}}
\caption{\textbf{Qualitative ablation results.} Full model produces most realistic smoke with proper semi-transparency and boundary softness.}
\label{fig:ablation_vis}
\end{figure}

\FloatBarrier

\subsection{Qualitative Comparison}

Figure~\ref{fig:qualitative} presents side-by-side visual comparisons across methods. \method~produces smoke with:
\begin{itemize}
    \item \textbf{Semi-transparent boundaries} that naturally blend with backgrounds (unlike GANs and vanilla diffusion)
    \item \textbf{Realistic color} (appropriate grays, browns, and blues depending on background)
    \item \textbf{Multi-scale structure} (fine wisps and large plumes coexist coherently)
    \item \textbf{Context-aware density} (thicker near source, thinner at edges)
\end{itemize}

Failure cases (bottom row) show occasional over-opacity or fog confusion, which we analyze in Section VII.

\begin{figure*}[t]
\centering
\IfFileExists{qualitative_comparison.pdf}{%
\includegraphics[width=0.98\textwidth]{qualitative_comparison.pdf}%
}{\PackageError{FireSmokeGenNet}{Missing qualitative_comparison.pdf}{Upload the experimentally generated qualitative comparison panel before submission}}
\caption{\textbf{Qualitative comparison.} Rows 1-3: Successful generations. Row 4: Failure cases (fog confusion, over-opacity). \method~produces most realistic boundaries and transparency.}
\label{fig:qualitative}
\end{figure*}

\section{Discussion}
\label{sec:discussion}

\subsection{Why Does \method~Work?}

Three principles explain our method's effectiveness:

\paragraph{P1: Decoupled Conditioning Prevents Interference.}
By processing the mask and masked image through separate ResNet-18 and ResNet-50 branches, respectively, the framework preserves modality-specific representations before fusion. The frozen ImageNet-pretrained encoders provide complementary structural and contextual features, while hierarchical injection supplies these features at the corresponding U-Net stages.

\paragraph{P2: Stochastic Boundary Regularization Matches Physical Reality.}
MRDL's morphological perturbations simulate the natural variability of smoke interfaces. By training the model to produce consistent denoising predictions under local boundary perturbations, MRDL reduces sensitivity to a single rigid annotation contour. This produces naturally diffuse gradients that align with real smoke statistics (KL divergence 0.041 vs. 0.187 for baseline).

\paragraph{P3: Hierarchical Injection Respects Semantics.}
Low-resolution layers receive global mask constraints (where should smoke appear?). High-resolution layers receive fine image details (what textures and colors?). This aligns with how smoke should be generated: first approximate shape, then refine texture. The JCA module at mid-resolutions enables cross-modal interaction where both structure and texture are important.

\subsection{Comparison with Prior Work}

\paragraph{vs. GAN-based Methods.}
GAN-based smoke synthesis \cite{rs12223715} can be affected by limited mode coverage and training instability. In the present experiments, the diffusion baselines produced more diverse smoke appearances and stronger downstream detection performance than the evaluated GAN baseline.

\paragraph{vs. Standard Diffusion Inpainting.}
Flat concatenation \cite{rombach2022highresolutionimagesynthesislatent, wang2024flamediffuserwildfireimage} forces the network to learn ad-hoc disentanglement of noise, mask, and image signals. Our dual-branch design with JCA provides explicit separation, leading to better boundary quality (8.1\% higher PSNR, 46.7\% better domain generalization under weather shift).

\paragraph{vs. Simulation-Based Methods.}
Physics-based rendering \cite{6053841} requires solving Navier-Stokes equations (computationally expensive) and manual parameter tuning. Our diffusion-based approach learns the distribution directly from data, requiring no physical simulation.

\subsection{Limitations and Failure Modes}
\label{sec:limitations}

Despite its performance relative to the evaluated baselines, \method~exhibits several limitations:

\paragraph{L1: Fog-Smoke Confusion.}
In high-haze conditions, generated smoke is observed in some generated
samples to resemble fog rather than distinct smoke plumes. This occurs
because the training data contains both fog and smoke examples, and the
model sometimes confuses them. \textbf{[PENDING VERIFICATION]} a specific
failure-rate percentage was previously reported here without a defined
denominator, evaluation set, annotation protocol, or confidence interval;
it has been removed pending a properly specified failure-rate study (a
fixed held-out evaluation set, an explicit annotation protocol for
labeling a generation as a fog-confusion failure, and a reported
confidence interval). Potential solution: incorporate explicit weather
conditioning or use contrastive learning to separate fog and smoke
representations.

\paragraph{L2: Over-Opacity in High-Contrast Scenes.}
When smoke should be extremely faint (early-stage fires, distant plumes),
the model is observed in some generated samples to produce overly opaque
outputs. This reflects a bias in the training data toward visible smoke.
\textbf{[PENDING VERIFICATION]} as with L1, a specific failure-rate
percentage was previously reported here without a defined denominator,
evaluation set, annotation protocol, or confidence interval; it has been
removed pending the same properly specified failure-rate study. Potential
solution: augment training with more examples of faint smoke or use
density prediction as an auxiliary task.

\paragraph{L3: Limited Quality-Annotation Set.}
The task-specific VLM is trained from 150 human-rated generated images. Although the observed agreement supports consistent annotation within this set, the limited sample size constrains conclusions about score calibration and transfer to other smoke datasets. Future work should use a larger independently held-out rating set and report rank correlation and prediction error.

\paragraph{L4: Computational Cost.}
With 50 DDIM sampling steps, \method~requires approximately 2 seconds per $512\times512$ image on an A100 GPU. While acceptable for offline dataset generation, this is too slow for real-time applications. Potential solution: distill to a few-step sampler using progressive distillation \cite{salimans2022progressivedistillationfastsampling} or consistency models \cite{song2023consistencymodels}.

\subsection{Future Directions}

\paragraph{F1: Video Smoke Synthesis.}
Extend the framework to latent video diffusion. This would enable explicit evaluation of temporal consistency and smoke dynamics, including plume movement, dissipation, and growth.

\paragraph{F2: Physical Grounding.}
Incorporate fluid dynamics priors into diffusion guidance. The score function could be augmented with a term that penalizes non-physical flow fields, ensuring generated smoke obeys conservation laws.

\paragraph{F3: Active Learning.}
Use detection model uncertainty to identify challenging scenes and target them for additional synthetic generation. This could create a virtuous cycle: detector $\rightarrow$ identifies weaknesses $\rightarrow$ generator synthesizes targeted examples $\rightarrow$ detector improves.

\paragraph{F4: Real-Time Adaptation.}
Develop lightweight adapters that can fine-tune the generator on the fly using detection feedback, enabling deployment-specific customization.

\paragraph{F5: Early-Detection / Time-to-Alarm Evaluation.}
The present experiments evaluate still-image smoke detection performance
(mAP@50, mAP@50--95, Precision, Recall) and do not measure quantities
specific to \emph{early} wildfire monitoring, such as detection delay,
time-to-alarm relative to fire ignition, sensitivity to early-stage
(faint, low-opacity) smoke specifically, video-level false-alarm rate, or
temporal stability of detections across consecutive frames. A dedicated
early-detection study -- e.g., replaying FIgLib/HPWREN video sequences
frame-by-frame and measuring the elapsed time between ignition and the
first sustained detection, both with and without \method-augmented
training data -- is required before any claim of improved \emph{early}
monitoring can be supported; this is reflected in the moderated title of
this manuscript (``Detection'' rather than ``Monitoring'') until such a
study is completed.

\section{Conclusion}
\label{sec:conclusion}

We presented \method, a comprehensive diffusion-based framework for wildfire smoke synthesis that addresses three critical limitations of prior work: flat conditioning, deterministic boundaries, and lack of quality control. Our contributions are:

\begin{enumerate}
    \item A hierarchical dual-branch conditioning mechanism with joint cross-attention that provides multi-scale spatial control while preventing feature interference.
    \item Mask Random Difference Loss (MRDL), a stochastic boundary regularization that models boundary uncertainty via morphological perturbations, achieving a boundary-softness KL divergence (0.041) that is the closest to the real-smoke reference among the evaluated methods.
    \item A multimodal quality filtering pipeline using fine-tuned Qwen2-VL that retains only the top 30\% of samples, improving downstream detection mAP@50 by 8.1 percentage points (a 10.8\% relative improvement, computed as $(0.829-0.748)/0.748\times100$) over random selection at the same retained dataset size; Section~\ref{sec:filtering_confound} reports the additional equal-size controls required to fully isolate this effect from the accompanying dataset-size change.
\end{enumerate}

Across the evaluated baselines, \method~achieves PSNR of 28.10, SSIM of 0.88, and YOLOv13 mAP@50 of 0.829, corresponding to a 21.4\% relative improvement over real-only training. Under challenging domain shifts, \method~achieves up to 46.7\% relative mAP@50 improvement, with all gains statistically significant ($p<0.01$). Ablation studies confirm the necessity of each component, with the full model outperforming variants lacking JCA or MRDL by 8.6\% and 7.0\% respectively.

These results indicate that task-aware diffusion synthesis can complement limited real data for wildfire-smoke monitoring. Our code, dataset, and pretrained models will be publicly released to facilitate further research.

\section*{Acknowledgment}

This work was supported in part by the National Science and Technology Council, Taiwan, under Grant NSTC 112-2221-E-A49-099-MY3. The authors thank Gama Innovation Inc. for providing computational resources and domain expertise. 
\bibliographystyle{IEEEtran}
\bibliography{ref}

% that's all folks
\end{document}
